\section{基于线性关系的分类}
\label{sec:classification-geometry}

第\ref{sec:classic-overview-linear}节的线性得分在分类中用于比较类别。本节从LDA、感知机和SVM考察三种学习依据：类统计量、分类错误反馈和间隔。相同的线性边界可以来自不同准则；随后用简介已建立的核方法扩展SVM的特征空间。

LDA兼有投影与概率解释，本节推导其共享协方差形式，第\ref{sec:classification-probability}节再从GDA模型族说明一般协方差设定。Fisher投影的统计量、感知机的可分性条件与软间隔SVM的误差容忍约束不同对象，不能排成假设强度依次递减的序列。

\subsection{线性判别分析（LDA）}
\label{sec:classification-lda}

如果用一个线性得分区分类别，怎样确定各特征的权重和分类阈值？一种办法是先描述每一类样本在空间中的分布，再比较新样本来自各类的可能性。\term{线性判别分析}（linear discriminant analysis, LDA）在高斯生成模型下采用这条路线：估计类均值、共享协方差和类先验，再由贝叶斯法则得到线性判别函数。

LDA还有一个几何起点。Fisher在1936年的鸢尾花分类研究中寻找一种线性组合，使两类均值的差异相对于类内波动尽可能大。这一准则不要求样本服从高斯分布；它与高斯生成模型可以给出相同的判别方向，但概率解释和适用条件有所不同\scite{fisher1936}

从两类的线性组合转向多个总体的判别，问题不再只是选择一条轴，还包括需要保留多少个方向。Rao在1948年系统讨论了多总体的生物分类与判别表示，这也为理解LDA如何从分类规则延伸为有监督降维方法提供了历史背景\scite{rao1948}

\subsubsection{模型结构}
\label{sec:classification-lda-model}

沿用训练集$D_n=\{(\bm{x}_i,y_i)\}_{i=1}^{n}$，其中$\bm{x}_i\in\mathbb{R}^{d}$、$y_i\in\{1,\ldots,C\}$。生成式LDA假设
\begin{equation}
  \label{eq:classification-lda-distribution}
  \begin{aligned}
    \mathbb{P}(y=c)&=\pi_c,\qquad \pi_c>0,\quad
      \sum_{c=1}^{C}\pi_c=1,\\
    p(\bm{x}\mid y=c)&=\mathcal{N}(\bm{x};\bm{\mu}_c,\bm{\Sigma}),
      \qquad \bm{\Sigma}\succ0 .
  \end{aligned}
\end{equation}
这里$\bm{\mu}_c\in\mathbb{R}^{d}$描述第$c$类的中心，$\bm{\Sigma}\in\mathbb{R}^{d\times d}$描述类内各方向的波动及特征之间的相关性。共享协方差意味着各类可以有不同中心，但在中心附近具有相同的椭球形状和朝向；并不要求各个特征相互独立。

在零一损失下，应选择后验概率最大的类别。取对数并展开高斯密度，有
\begin{equation}
  \label{eq:classification-lda-score}
  \begin{aligned}
    \log\{\pi_c p(\bm{x}\mid y=c)\}
      &=a(\bm{x})+\delta_c(\bm{x}),\\
    \delta_c(\bm{x})
      &=\bm{x}^{\top}\bm{\Sigma}^{-1}\bm{\mu}_c\\
      &\quad-\frac12\bm{\mu}_c^{\top}\bm{\Sigma}^{-1}\bm{\mu}_c
        +\log\pi_c .
  \end{aligned}
\end{equation}
其中$a(\bm{x})$收集了$-\bm{x}^{\top}\bm{\Sigma}^{-1}\bm{x}/2$、对数行列式和高斯归一化常数，均与类别无关。正是共享协方差使二次项在类别比较时抵消，留下关于$\bm{x}$的仿射函数。任意两类得分相等的位置通常为超平面；多分类时，各类别区域由若干这样的比较共同确定，而不是全体类别只共用一个分割面。

训练后把估计参数代入$\delta_c$，得到$\hat{\delta}_c$。预测规则及相应的模型后验为
\begin{equation}
  \label{eq:classification-lda-prediction}
  \begin{aligned}
    \hat f_{\mathrm{LDA}}(\bm{x})
      &=\argmax_{c\in\{1,\ldots,C\}}\hat{\delta}_c(\bm{x}),\\
    \hat p_c(\bm{x})
      &=\frac{\exp\{\hat{\delta}_c(\bm{x})\}}
        {\sum_{j=1}^{C}\exp\{\hat{\delta}_j(\bm{x})\}} .
  \end{aligned}
\end{equation}
本节对$\argmax$的平局统一选择索引最小的类别。若放宽共享协方差约束，让每类使用自己的$\bm{\Sigma}_c$，二次项便不再普遍抵消，得到\term{二次判别分析}（quadratic discriminant analysis, QDA），即允许二次决策边界的高斯分类器\scite{hastie2009}

\subsubsection{参数学习}
\label{sec:classification-lda-learning}

\paragraph{生成模型的最大似然估计}
有标签样本已经给出了每个观测所属的高斯分量，因此这里不需要推断类别归属。令$n_c=\sum_{i=1}^{n}\mathbb{1}\{y_i=c\}$，暂设各类都有样本。忽略与参数无关的常数，对数似然为
\begin{equation}
  \label{eq:classification-lda-likelihood}
  \begin{aligned}
    \ell
      &=\sum_{c=1}^{C}n_c\log\pi_c
        -\frac n2\log\det\bm{\Sigma}\\
      &\quad-\frac12\sum_{c=1}^{C}\sum_{i:y_i=c}
        (\bm{x}_i-\bm{\mu}_c)^{\top}
        \bm{\Sigma}^{-1}(\bm{x}_i-\bm{\mu}_c).
  \end{aligned}
\end{equation}
对先验加入约束$\sum_c\pi_c=1$，拉格朗日驻点给出$n_c/\pi_c$为同一常数；归一化后该常数为$n$。对$\bm{\mu}_c$求导，则有$\bm{\Sigma}^{-1}\sum_{i:y_i=c}(\bm{x}_i-\bm{\mu}_c)=\bm{0}$。由此得到
\begin{equation}
  \label{eq:classification-lda-mean-prior}
  \hat\pi_c=\frac{n_c}{n},\qquad
  \hat{\bm{\mu}}_c=\frac1{n_c}\sum_{i:y_i=c}\bm{x}_i .
\end{equation}
这两个估计分别是类别出现频率和该类的样本均值。

把均值代回似然，定义\term{类内散度矩阵}（within-class scatter matrix）
\begin{equation}
  \label{eq:classification-lda-within}
  \bm{S}_{\mathrm{W}}
    =\sum_{c=1}^{C}\sum_{i:y_i=c}
      (\bm{x}_i-\hat{\bm{\mu}}_c)
      (\bm{x}_i-\hat{\bm{\mu}}_c)^{\top}.
\end{equation}
它累加每个样本相对于本类中心的偏离。用精度矩阵$\bm{\Omega}=\bm{\Sigma}^{-1}$改写剩余目标，得到$n\log\det\bm{\Omega}/2-\operatorname{tr}(\bm{\Omega}\bm{S}_{\mathrm{W}})/2$。其驻点满足$n\bm{\Omega}^{-1}=\bm{S}_{\mathrm{W}}$，故当$\bm{S}_{\mathrm{W}}\succ0$时，
\begin{equation}
  \label{eq:classification-lda-covariance}
  \hat{\bm{\Sigma}}=\frac1n\bm{S}_{\mathrm{W}} .
\end{equation}
对$\bm{\Omega}$的目标是严格凹的，因此该驻点给出最大似然解。

式~\eqref{eq:classification-lda-covariance}的分母是$n$，不是$n-C$。这种合并各类类内变动的估计称为\term{合并协方差}（pooled covariance），其无偏形式为
\[
  \hat{\bm{\Sigma}}_{\mathrm{U}}
    =\frac{\bm{S}_{\mathrm{W}}}{n-C},\qquad n>C .
\]
在共享协方差假设下，给定各类样本数，有$\mathbb{E}[\bm{S}_{\mathrm{W}}\mid y_1,\ldots,y_n]=(n-C)\bm{\Sigma}$，因为估计每个类中心各消耗一个自由度。最大似然与无偏估计是不同准则；两种估计渐近相同，但有限样本下有所区别。特别是先验不等时，整体缩放协方差会改变几何得分相对于$\log\hat\pi_c$的权重，预测边界未必相同。

闭式估计不等于没有计算代价。对稠密数据，累积类内散度需要$O(nd^2)$运算，矩阵分解需要$O(d^3)$运算。实现时分解一次$\hat{\bm{\Sigma}}$，再解$C$个线性方程组$\hat{\bm{\Sigma}}\bm{a}_c=\hat{\bm{\mu}}_c$，预存$\bm{a}_c$和常数得分项；不必显式求逆，预测一个样本只需$O(Cd)$运算。若$\bm{S}_{\mathrm{W}}$奇异，正定参数域内上述最大似然解并不存在，需要降低维数或作正则化，不能直接使用逆矩阵公式。

\paragraph{Fisher投影的学习}
生成模型之外，还可以直接比较投影后的类间分离与类内波动。令总样本均值$\hat{\bm{\mu}}=\sum_c(n_c/n)\hat{\bm{\mu}}_c$，定义\term{类间散度矩阵}（between-class scatter matrix）
\[
  \bm{S}_{\mathrm{B}}
    =\sum_{c=1}^{C}n_c
      (\hat{\bm{\mu}}_c-\hat{\bm{\mu}})
      (\hat{\bm{\mu}}_c-\hat{\bm{\mu}})^{\top}.
\]
它衡量各类中心相对于总中心的分散程度。投影到方向$\bm{w}$后，两种散度分别变为$\bm{w}^{\top}\bm{S}_{\mathrm{B}}\bm{w}$和$\bm{w}^{\top}\bm{S}_{\mathrm{W}}\bm{w}$。\term{Fisher判别准则}（Fisher discriminant criterion）取两者之比。以下假设$\bm{S}_{\mathrm{W}}\succ0$，使每个非零方向的分母都为正：
\begin{equation}
  \label{eq:classification-fisher-criterion}
  J(\bm{w})
    =\frac{\bm{w}^{\top}\bm{S}_{\mathrm{B}}\bm{w}}
      {\bm{w}^{\top}\bm{S}_{\mathrm{W}}\bm{w}},
    \qquad \bm{w}\ne\bm{0}.
\end{equation}
分子鼓励类中心远离，分母抑制同类样本投影后过于分散。

图~\ref{fig:classification-lda-projection}固定同一组二维教学样本，只改变投影方向。沿$\bm{w}_{\mathrm{A}}=(1,0)^{\top}$，两类投影均值相差$2$，而各类仅在宽度$0.6$的区间内波动；沿$\bm{w}_{\mathrm{B}}=(3/5,4/5)^{\top}$，均值差缩为$1.2$，同类投影却散布得更宽，异类点在数轴上交错。由图注中的构造可算得$\bm{S}_{\mathrm{W}}=\operatorname{diag}(2/5,8)$、$\bm{S}_{\mathrm{B}}=\operatorname{diag}(8,0)$，所以$J(\bm{w}_{\mathrm{A}})=20$，$J(\bm{w}_{\mathrm{B}})=180/329$。这个比较衡量的是同一批样本投影后的相对分离程度，不是测试错误率。

\begin{figure}[htbp]
  \centering
  % 两类中心为 (-1,0)、(1,0)，共用四个精确十进制偏移；不使用随机点或抖动。
% 两面板均由同一偏移表生成，投影坐标直接计算为 wx*x1+wy*x2。
\begin{tikzpicture}[
  x=0.1\linewidth,y=0.1\linewidth,
  font=\footnotesize,
  every node/.style={inner sep=1pt}
]
  \path[use as bounding box] (0,-4.3) rectangle (10,2.4);
  \foreach \shift/\id/\wx/\wy/\pattern/\comparison/\ratio in {
    2.3/A/1/0/solid/较大/20,
    7.6/B/0.6/0.8/dashed/较小/{180/329}
  } {
    \begin{scope}[shift={(\shift,0)}]
      \node[font=\heiti\footnotesize] at (0,2.1)
        {（\id）相对分离\comparison};
      \draw[BookMuted,thin,-{Stealth[length=1.3mm]}]
        (-1.75,0) -- (1.85,0) node[right] {$x_1$};
      \draw[BookMuted,thin,-{Stealth[length=1.3mm]}]
        (0,-1.25) -- (0,1.4) node[above] {$x_2$};
      \foreach \tick in {-1,1} {
        \draw[BookMuted] (\tick,-0.04) -- (\tick,0.04);
        \node[below] at (\tick,-0.06) {$\tick$};
        \draw[BookMuted] (-0.04,\tick) -- (0.04,\tick);
        \node[left] at (-0.06,\tick) {$\tick$};
      }
      \node[below left] at (0,0) {$0$};
      \draw[BookTeal,thick,\pattern,-{Stealth[length=1.5mm]}]
        (0,0) -- ({1.45*\wx},{1.45*\wy})
        node[above right] {$\bm{w}_{\mathrm{\id}}$};

      \node at (0,-1.65) {$z_{\mathrm{\id}}=\bm{w}_{\mathrm{\id}}^{\top}\bm{x}$};
      \foreach \class/\cx/\colour/\row in {
        1/-1/BookBlue/-2.1,
        2/1/BookAmber/-2.5
      } {
        \draw[BookRule] (-1.75,\row) -- (1.75,\row);
        \foreach \dx/\dy in {-0.3/-1,-0.1/1,0.1/1,0.3/-1} {
          \pgfmathsetmacro{\px}{\cx+\dx}
          \pgfmathsetmacro{\pz}{\wx*\px+\wy*\dy}
          \ifnum\class=1
            \fill[\colour] (\px,\dy) circle (1.5pt);
            \fill[\colour] (\pz,\row) circle (1.5pt);
          \else
            \fill[\colour] ([xshift=-1.5pt,yshift=-1.5pt]\px,\dy)
              rectangle ++(3pt,3pt);
            \fill[\colour] ([xshift=-1.5pt,yshift=-1.5pt]\pz,\row)
              rectangle ++(3pt,3pt);
          \fi
        }
        \pgfmathsetmacro{\mz}{\wx*\cx}
        \draw[\colour,line width=0.8pt]
          (\mz,{\row-0.13}) -- (\mz,{\row+0.13});
      }
      \draw[BookMuted,thin,-{Stealth[length=1.3mm]}]
        (-1.75,-2.9) -- (1.85,-2.9) node[right] {$z$};
      \foreach \tick in {-1,0,1} {
        \draw[BookMuted] (\tick,-2.94) -- (\tick,-2.86);
        \node[below] at (\tick,-2.97) {$\tick$};
      }
      \node at (0,-3.5) {$J(\bm{w}_{\mathrm{\id}})=\ratio$};
    \end{scope}
  }
  \fill[BookBlue] (1.8,-4.05) circle (1.5pt);
  \node[right] at (2,-4.05) {类1};
  \fill[BookAmber] ([xshift=-1.5pt,yshift=-1.5pt]4,-4.05)
    rectangle ++(3pt,3pt);
  \node[right] at (4.2,-4.05) {类2};
  \draw[BookInk,line width=0.8pt] (6.1,-4.18) -- (6.1,-3.92);
  \node[right] at (6.3,-4.05) {投影均值};
\end{tikzpicture}

  \caption{同一点集在两个单位方向上的投影。两类各有4点，中心分别为$(-1,0)^{\top}$、$(1,0)^{\top}$，都加上偏移$(-0.3,-1)^{\top}$、$(-0.1,1)^{\top}$、$(0.1,1)^{\top}$、$(0.3,-1)^{\top}$；这是人为构造，坐标无量纲。两行只改变黄色投影方向：$\bm{w}_{\mathrm{A}}=(1,0)^{\top}$与$\bm{w}_{\mathrm{B}}=(3/5,4/5)^{\top}$。右侧按$z=\bm{w}^{\top}\bm{x}$把蓝圆和红方放到同一数轴，彩色短竖线标出各类投影均值；两行共用数轴尺度。}
  \label{fig:classification-lda-projection}
\end{figure}

由于$J$不受$\bm{w}$的非零倍数缩放影响，可以约束$\bm{w}^{\top}\bm{S}_{\mathrm{W}}\bm{w}=1$。对$\bm{w}^{\top}\bm{S}_{\mathrm{B}}\bm{w}-\lambda(\bm{w}^{\top}\bm{S}_{\mathrm{W}}\bm{w}-1)$求导，得到\term{广义特征值问题}（generalized eigenvalue problem）：
\begin{equation}
  \label{eq:classification-fisher-eigenproblem}
  \bm{S}_{\mathrm{B}}\bm{w}
    =\lambda\bm{S}_{\mathrm{W}}\bm{w}.
\end{equation}
驻点还需比较大小。令$\bm{v}=\bm{S}_{\mathrm{W}}^{1/2}\bm{w}$，则
\[
  J(\bm{w})=\frac{\bm{v}^{\top}\bm{A}\bm{v}}{\bm{v}^{\top}\bm{v}},
  \qquad
  \bm{A}=\bm{S}_{\mathrm{W}}^{-1/2}
    \bm{S}_{\mathrm{B}}\bm{S}_{\mathrm{W}}^{-1/2}.
\]
把$\bm{v}$展开到对称半正定矩阵$\bm{A}$的正交特征向量上，上式就是各特征值的加权平均，权重非负且总和为$1$。最大值因而是最大特征值；取对应特征向量，再乘$\bm{S}_{\mathrm{W}}^{-1/2}$，便得到最优Fisher方向。若最大特征值有重数，最优方向未必唯一。

二分类时，令$\bm{h}=\hat{\bm{\mu}}_1-\hat{\bm{\mu}}_2\ne\bm{0}$，由总均值的定义可得
\[
  \bm{S}_{\mathrm{B}}=\frac{n_1n_2}{n}\bm{h}\bm{h}^{\top},
  \qquad
  \bm{w}_{\mathrm{F}}\propto\bm{S}_{\mathrm{W}}^{-1}\bm{h}.
\]
生成式LDA的法向量则为$\bm{w}_{\mathrm{G}}=\hat{\bm{\Sigma}}^{-1}\bm{h}=n\bm{S}_{\mathrm{W}}^{-1}\bm{h}$，两者方向相同。但完整的生成式分类规则还比较
\begin{equation}
  \label{eq:classification-lda-binary-threshold}
  \hat\delta_1(\bm{x})-\hat\delta_2(\bm{x})
    =\bm{w}_{\mathrm{G}}^{\top}
      \left(\bm{x}-\frac{\hat{\bm{\mu}}_1+\hat{\bm{\mu}}_2}{2}\right)
      +\log\frac{\hat\pi_1}{\hat\pi_2}
\end{equation}
的正负。Fisher准则本身只给方向，不给类先验、后验概率或这一阈值。两类均值相同时，均值散度准则没有正特征值方向，而生成模型仍可按先验选择类别。由此可见，非高斯数据上仍能使用Fisher投影，但不能据此沿用高斯LDA的贝叶斯最优性结论。

\subsubsection{理论分析}
\label{sec:classification-lda-analysis}

\paragraph{收敛性分析}
生成式LDA的均值、先验与协方差都有闭式估计，因此这里考察的是样本量增加时估计是否趋近真实参数，而不是某个迭代过程何时停止。\term{一致性}（consistency）正是这种渐近性质；它不表示某个有限训练集上的估计已经准确。

\begin{theorem}[LDA参数估计的一致性]
  \label{thm:classification-lda-consistency}
  设$d$和$C$固定，训练样本独立同分布，真实分布满足式~\eqref{eq:classification-lda-distribution}，且所有$\pi_c^\ast>0$、$\bm{\Sigma}^\ast\succ0$。当$n\to\infty$时，
  \[
    \hat\pi_c\xrightarrow{\mathbb{P}}\pi_c^\ast,\qquad
    \hat{\bm{\mu}}_c\xrightarrow{\mathbb{P}}\bm{\mu}_c^\ast,\qquad
    \hat{\bm{\Sigma}}\xrightarrow{\mathbb{P}}\bm{\Sigma}^\ast .
  \]
  所有类别非空且$\hat{\bm{\Sigma}}$正定的概率趋于$1$。
\end{theorem}

\begin{proof}
  为处理随机的$n_c$，把类均值写成两个全样本平均之比：
  \[
    \hat{\bm{\mu}}_c
      =\frac{n^{-1}\sum_{i=1}^{n}\bm{x}_i\mathbb{1}\{y_i=c\}}
        {n^{-1}\sum_{i=1}^{n}\mathbb{1}\{y_i=c\}} .
  \]
  大数定律使分母收敛到$\pi_c^\ast>0$，分子收敛到$\pi_c^\ast\bm{\mu}_c^\ast$，因而先验和均值都一致。尚未观察到某一类时可以任意约定其均值；该事件的概率趋于零，不影响结论。

  共享协方差估计可展开为
  \[
    \hat{\bm{\Sigma}}
      =\frac1n\sum_{i=1}^{n}\bm{x}_i\bm{x}_i^{\top}
        -\sum_{c=1}^{C}\hat\pi_c
          \hat{\bm{\mu}}_c\hat{\bm{\mu}}_c^{\top}.
  \]
  高斯分布有有限二阶矩，故第一项收敛到
  \[
    \mathbb{E}[\bm{x}\bm{x}^{\top}]
      =\bm{\Sigma}^\ast+
        \sum_{c=1}^{C}\pi_c^\ast
          \bm{\mu}_c^\ast(\bm{\mu}_c^\ast)^{\top}.
  \]
  第二项消去类中心对二阶矩的贡献，余下$\bm{\Sigma}^\ast$。固定维数下，矩阵逐元素收敛也给出算子范数收敛；由$\bm{\Sigma}^\ast$的最小特征值严格为正，估计矩阵以趋于$1$的概率可逆。上述步骤实际上可由强大数定律加强为几乎必然收敛。
\end{proof}

矩阵求逆和$\log\pi_c$在这些条件下连续，所以在参数几乎必然收敛的事件上，对每个固定输入$\bm{x}$，有$\hat\delta_c(\bm{x})\to\delta_c^\ast(\bm{x})$。但$\argmax$在平局处不连续：在真实最大得分唯一的输入处，估计类别最终等于该输入的唯一贝叶斯类别；平局处则未必如此。若最优类别平局的输入集合概率为零，就能进一步推出随机输入上的预测标签收敛。不同贝叶斯分类器在平局处可以选择不同标签，却有相同风险。

\paragraph{泛化分析}
参数趋近真实值以后，还需要说明这如何影响独立测试样本上的错误率。分类风险的一致性不必要求标签处处收敛。记$p_c^\ast(\bm{x})$为真实后验，$f^\ast$任选一个贝叶斯规则。由于$\hat f_{\mathrm{LDA}}$最大化$\hat p_c$，逐点有
\begin{equation}
  \label{eq:classification-lda-excess-risk}
  \begin{aligned}
    0&\leq p_{f^\ast(\bm{x})}^\ast(\bm{x})
      -p_{\hat f_{\mathrm{LDA}}(\bm{x})}^\ast(\bm{x})\\
      &\leq 2\max_{1\leq c\leq C}
          |\hat p_c(\bm{x})-p_c^\ast(\bm{x})| .
  \end{aligned}
\end{equation}
这个上界来自在两项之间分别加减相应的估计后验，并利用$\hat p_{\hat f_{\mathrm{LDA}}}\geq\hat p_{f^\ast}$。在参数几乎必然收敛的事件上，后验估计也对每个输入收敛；右端趋于零且不超过$2$。对与训练集独立、来自同一分布的测试输入积分，再用支配收敛定理，得到
\[
  R(\hat f_{\mathrm{LDA}}\mid D_n)
    \longrightarrow R^\ast
    \quad\text{几乎必然},
\]
其中$R(\hat f_{\mathrm{LDA}}\mid D_n)$是给定训练集后的测试错误率，$R^\ast$是贝叶斯错误率。有限样本中无法拟合时可暂用任意固定分类器，不改变这一渐近结论。这是模型正确指定、固定维数下的保证，不是对任意分布或$d$随$n$快速增长的普遍保证，也未给出有限样本的收敛速率。Fisher准则在当前训练集上取得最大值，不能替代这一风险分析。

\subsubsection{支持有监督降维}
\label{sec:classification-lda-reduction}

如果类别多于两个，一条投影轴可能不足以保留类中心之间的差异。把多个广义特征向量作为$\bm{W}\in\mathbb{R}^{d\times r}$的列，就可以把输入转换为$\bm{z}=\bm{W}^{\top}(\bm{x}-\hat{\bm{\mu}})\in\mathbb{R}^{r}$。以下仍假设$\bm{S}_{\mathrm{W}}\succ0$，并取$1\leq r\leq d$。

为避免重复选择同一方向，可要求不同投影轴在类内散度度量下正交：
\begin{equation}
  \label{eq:classification-lda-subspace}
  \max_{\bm{W}}\operatorname{tr}
       (\bm{W}^{\top}\bm{S}_{\mathrm{B}}\bm{W})
  \quad\text{满足}\quad
  \bm{W}^{\top}\bm{S}_{\mathrm{W}}\bm{W}=\bm{I}_r .
\end{equation}
在前述白化坐标中，该问题选择$\bm{A}$最大的$r$个特征值对应的正交方向，因而也就是式~\eqref{eq:classification-fisher-eigenproblem}中最大的$r$个广义特征值方向。

为什么有用的方向至多为$C-1$个？令$\bm{M}\in\mathbb{R}^{d\times C}$的第$c$列为$\sqrt{n_c}(\hat{\bm{\mu}}_c-\hat{\bm{\mu}})$，则
\[
  \bm{S}_{\mathrm{B}}=\bm{M}\bm{M}^{\top},\qquad
  \bm{M}
  \begin{pmatrix}\sqrt{n_1}\\ \vdots\\ \sqrt{n_C}\end{pmatrix}
  =\bm{0}.
\]
右侧列向量非零，说明$\bm{M}$的$C$列线性相关。因此$\operatorname{rank}(\bm{S}_{\mathrm{B}})\leq\min(d,C-1)$；可逆白化不改变秩，$\bm{A}$也至多有这么多个正特征值。若类中心还有额外的线性依赖，实际维数更低。这个限制针对类均值散度，不表示任意分类任务的全部信息都能无损压缩到$C-1$维；例如均值相同而协方差不同的类别，可能被该准则忽略。

LDA投影使用了标签，因此训练和验证时必须把它作为学习流程的一部分，在每个训练划分内重新估计投影，不能先用全体样本标签降维再评估。得到的低维特征可以交给近邻分类器等后续模型。第\ref{chap:dimensionality-reduction}章按保留信息的目标比较降维方法；其中第\ref{sec:dim-pca}节的PCA保留输入方差而不使用类别标签，LDA则保留类别均值相对于类内波动的判别信息。Fisherfaces人脸识别方法就是借助Fisher准则构造低维判别表示的一个具体例子\scite{belhumeur1997}

\subsubsection{小结}
\label{sec:classification-lda-summary}

LDA给线性得分提供了两种解释：共享协方差高斯模型给出概率、方向和阈值，Fisher准则则直接选择类中心相对分离的投影。两者在二分类且采用同一类内散度时给出相同方向，但概率解释和风险保证仍依赖额外假设。作为特征变换，LDA也用于语音识别；例如Haeb-Umbach与Ney研究了按语音类别学习线性变换，再与后续识别模型结合的方案\scite{haebumbach1992}

选用LDA时，需要同时考察分布近似和协方差的可估计性。偏离高斯不必然使分类失效，却使前述最优性证明不再适用；异常样本也会影响均值与协方差。由于$\operatorname{rank}(\bm{S}_{\mathrm{W}})\leq\min(d,n-C)$，当$n-C<d$时估计必然奇异，即使样本稍多于这个界也可能不稳定。降维或收缩协方差可以改善可计算性，但引入了新的选择及偏差，需要用独立验证评估。

\subsection{感知机}
\label{sec:classification-perceptron}

如果不估计类条件分布，能否直接根据分类结果调整边界？\term{感知机}（perceptron）的做法是用线性得分预测类别，遇到错误就修正权重。它不要求高斯分布；线性可分性是其有限次更新保证的条件，而不是运行算法的前提。Rosenblatt在1958年的论文中研究了感知机的学习机制，本节介绍的是其中广泛使用的线性阈值模型与误差修正规则\scite{rosenblatt1958}

这条研究路线也有硬件实现。康奈尔航空实验室建造的Mark I Perceptron将感知机思想实现为电子装置，但历史上的系统还包含感觉、关联和响应等组织，不能把它的全部结构等同于今天的单个线性分类器\scite{smithsonianMarkI}

\subsubsection{模型结构}
\label{sec:classification-perceptron-model}

本小节先使用二分类标签$y_i\in\{-1,+1\}$。为了使偏置在更新规则和收敛证明中始终处于同一坐标系，引入\term{增广向量}（augmented vector），即在输入末尾补一个恒为$1$的分量，并把偏置并入参数：
\begin{equation}
  \label{eq:classification-perceptron-augmented}
  \tilde{\bm{x}}=
    \begin{pmatrix}\bm{x}\\1\end{pmatrix},\qquad
  \tilde{\bm{w}}=
    \begin{pmatrix}\bm{w}\\b\end{pmatrix}
    \in\mathbb{R}^{d+1}.
\end{equation}
其中$\bm{w}\in\mathbb{R}^{d}$，$b\in\mathbb{R}$。分类规则为
\begin{equation}
  \label{eq:classification-perceptron-prediction}
  f(\bm{x})=\operatorname{sign}
    (\tilde{\bm{w}}^{\top}\tilde{\bm{x}})
    =\operatorname{sign}(\bm{w}^{\top}\bm{x}+b).
\end{equation}
约定$\operatorname{sign}(s)=+1$当$s\geq0$，否则为$-1$。非退化边界$\bm{w}^{\top}\bm{x}+b=0$是在$d$维输入空间中的$(d-1)$维超平面。它与二分类LDA具有相同的边界形式，但其参数不由类分布及先验推出，也不自动给出类别概率。

在这一表示下，正间隔可分是指存在单位向量$\bm{u}^\ast\in\mathbb{R}^{d+1}$和$\gamma>0$，使
\begin{equation}
  \label{eq:classification-perceptron-margin}
  \begin{aligned}
    y_i\langle\bm{u}^\ast,\tilde{\bm{x}}_i\rangle
      &\geq\gamma,\qquad i=1,\ldots,n,\\
    \|\bm{u}^\ast\|_2&=1,\qquad
    \|\tilde{\bm{x}}_i\|_2\leq R .
  \end{aligned}
\end{equation}
这里$R$约束的是增广输入，因此$R\geq\max_i\sqrt{\|\bm{x}_i\|_2^2+1}$。$\gamma$也是增广空间中的间隔；若$\bm{u}^\ast=(\bm{v}^{\top},a)^{\top}/\sqrt{\|\bm{v}\|_2^2+a^2}$，其分母包含偏置$a$，不同于原输入空间中以$\|\bm{v}\|_2$归一化的几何距离。有限训练集只要能被仿射超平面严格分开，就存在这样的正$\gamma$。

线性表示的边界也很具体。例如异或任务把$(0,0)$、$(1,1)$标为负类，把$(1,0)$、$(0,1)$标为正类。按上述平局规则，前两点要求$b<0$和$w_1+w_2+b<0$，后两点要求$w_1+b\geq0$和$w_2+b\geq0$。前两式相加与后两式相加，对同一个$w_1+w_2+2b$给出矛盾要求。Minsky与Papert在1969年的《Perceptrons》中系统研究了此类表示能力限制；改变输入特征或加入非线性隐藏层可以越过单个线性边界的限制，但那已经改变了模型族\scite{minsky1969}

\subsubsection{参数学习}
\label{sec:classification-perceptron-learning}

从$\tilde{\bm{w}}^{(0)}=\bm{0}$出发，若当前样本满足$y_i\langle\tilde{\bm{w}}^{(t)},\tilde{\bm{x}}_i\rangle\leq0$，便进行一次更新：
\begin{equation}
  \label{eq:classification-perceptron-update}
  \tilde{\bm{w}}^{(t+1)}
    =\tilde{\bm{w}}^{(t)}+\eta y_i\tilde{\bm{x}}_i,
    \qquad \eta>0 .
\end{equation}
这里$t$只数实际更新，学习率$\eta$保持不变。拆开最后一个分量，就是$\bm{w}\leftarrow\bm{w}+\eta y_i\bm{x}_i$和$b\leftarrow b+\eta y_i$。同一个样本的有符号得分增加$\eta\|\tilde{\bm{x}}_i\|_2^2$，所以更新确实朝修正当前判断的方向移动；但它可能同时破坏其他样本的正确分类，因而需要反复遍历。零得分也触发更新，即使按平局规则恰好预测为正确的正类，仍要求它离开边界。

\begin{algorithm}[带遍历上限的感知机训练]
  \label{alg:classification-perceptron}
  \begin{algorithmic}[1]
    \Require $D_n$，固定学习率$\eta>0$，最大遍历轮数$E\geq1$
    \Ensure 参数$\tilde{\bm{w}}$与是否已验证收敛的状态
    \State $\tilde{\bm{w}}\gets\bm{0}$
    \For{$e=1,\ldots,E$}
      \State $m\gets0$
      \For{$i=1,\ldots,n$}
        \State $\tilde{\bm{x}}_i\gets(\bm{x}_i^{\top},1)^{\top}$
        \If{$y_i\langle\tilde{\bm{w}},\tilde{\bm{x}}_i\rangle\leq0$}
          \State $\tilde{\bm{w}}\gets\tilde{\bm{w}}+\eta y_i\tilde{\bm{x}}_i$
          \State $m\gets m+1$
        \EndIf
      \EndFor
      \If{$m=0$}
        \State \Return $(\tilde{\bm{w}},\text{已收敛})$
      \EndIf
    \EndFor
    \State \Return $(\tilde{\bm{w}},\text{达到轮数上限})$
  \end{algorithmic}
\end{algorithm}

一整轮没有更新时，所有样本都是在同一个参数下得到严格正的有符号得分，才构成收敛证据。达到轮数上限只表示预算用尽，不能据此断言数据不可分。稠密实现每轮需要$O(nd)$运算，除训练集外只需$O(d)$参数空间；样本访问顺序可以改变输出的分隔面。

从优化角度看，定义\term{感知机损失}（perceptron loss）
\begin{equation}
  \label{eq:classification-perceptron-loss}
  \ell_{\mathrm{P}}(\tilde{\bm{w}};\bm{x}_i,y_i)
    =\max\{0,-y_i\langle\tilde{\bm{w}},\tilde{\bm{x}}_i\rangle\}.
\end{equation}
负得分处的梯度为$-y_i\tilde{\bm{x}}_i$，正得分处为$\bm{0}$；零得分处选择次梯度$-y_i\tilde{\bm{x}}_i$，就得到式~\eqref{eq:classification-perceptron-update}。这是逐样本次梯度更新；随机抽取样本时才相应成为随机次梯度下降。这里使用的是零间隔损失，不是标准的\term{合页损失}（hinge loss）$\max\{0,1-y_i\langle\tilde{\bm{w}},\tilde{\bm{x}}_i\rangle\}$。前者在零参数处就对所有样本取零，因此“感知机损失为零”本身不能证明分类正确；算法在零得分处仍更新的约定不可省略。

\subsubsection{理论分析}
\label{sec:classification-perceptron-analysis}

\paragraph{收敛性分析}
每次修正可能破坏已有判断，为什么这些修正不会无穷持续？关键是存在一个始终把所有样本分对的参考方向：更新朝该方向取得线性累积的进展，而参数范数的平方只能线性增长。Novikoff在1962年报告的证明清楚展示了这两个增长速度之间的矛盾\scite{novikoff1962}

\begin{theorem}[感知机的错误更新界]
  \label{thm:classification-perceptron-mistakes}
  若$D_n$满足式~\eqref{eq:classification-perceptron-margin}，从零参数出发，以固定$\eta>0$按式~\eqref{eq:classification-perceptron-update}更新，则对任意样本访问顺序，错误更新次数$M$满足
  \[
    M\leq\left(\frac R\gamma\right)^2 .
  \]
  这里错误更新包括零得分时的更新。若持续完整遍历训练集，且不因预算提前退出，算法在有限轮后得到所有训练样本上有符号得分均严格为正的参数。
\end{theorem}

\begin{proof}
  设第$t+1$次更新使用样本$i_t$。沿单位参考方向$\bm{u}^\ast$考察更新，得到
  \begin{align*}
    \langle\tilde{\bm{w}}^{(t+1)},\bm{u}^\ast\rangle
      &=\langle\tilde{\bm{w}}^{(t)},\bm{u}^\ast\rangle
        +\eta y_{i_t}\langle\tilde{\bm{x}}_{i_t},\bm{u}^\ast\rangle\\
      &\geq\langle\tilde{\bm{w}}^{(t)},\bm{u}^\ast\rangle+\eta\gamma .
  \end{align*}
  从零参数累加$M$次，便有
  \[
    \langle\tilde{\bm{w}}^{(M)},\bm{u}^\ast\rangle
      \geq M\eta\gamma .
  \]
  另一方面，展开更新后的平方范数：
  \begin{align*}
    \|\tilde{\bm{w}}^{(t+1)}\|_2^2
      &=\|\tilde{\bm{w}}^{(t)}\|_2^2
        +2\eta y_{i_t}
          \langle\tilde{\bm{w}}^{(t)},\tilde{\bm{x}}_{i_t}\rangle\\
      &\quad+\eta^2\|\tilde{\bm{x}}_{i_t}\|_2^2\\
      &\leq\|\tilde{\bm{w}}^{(t)}\|_2^2+\eta^2R^2 .
  \end{align*}
  不等式使用了更新条件：交叉项非正；并使用同一增广空间中的长度界$R$。再次累加，得到$\|\tilde{\bm{w}}^{(M)}\|_2^2\leq M\eta^2R^2$。由Cauchy--Schwarz不等式和$\|\bm{u}^\ast\|_2=1$，
  \[
    M\eta\gamma
      \leq\langle\tilde{\bm{w}}^{(M)},\bm{u}^\ast\rangle
      \leq\|\tilde{\bm{w}}^{(M)}\|_2
      \leq\eta R\sqrt{M}.
  \]
  $M=0$时结论直接成立；$M>0$时除以$\eta\sqrt M$再平方，得到所述界。因此更新不可能无限发生。持续完整遍历会在最后一次更新之后再检查所有样本；若仍有非正得分就还会更新，产生矛盾，故最终必有一整轮没有更新。
\end{proof}

这个界不显式含$n$，也不依赖访问顺序，但$R$和$\gamma$随数据而定；增加样本可能缩小可分间隔。它只约束更新次数，不包括寻找错误时检查正确样本的开销。对于算法~\ref{alg:classification-perceptron}的逐轮扫描，若不提前截断，至多需要$M+1$轮，包括最终的无更新确认轮。

若数据不可分，有限更新保证失效，持续遍历可能反复更新，此时可以限制训练预算并保留历史较好解，或采用允许间隔违约的在线方法。这些处理改变了停止或选解方式，仍须分别分析。

\paragraph{泛化分析}
错误更新界不是最终分类器的测试风险界。前述证明既没有独立同分布假设，也没有说明未来数据如何产生；仅凭某个训练集上的零误差，不能控制测试误差。研究问题从“修正会不会停止”转向“怎样从训练历史中选出风险小的分类器”后，还需要抽样条件和输出规则。Cesa-Bianchi、Conconi与Gentile在2004年系统分析了这种联系：利用独立同分布样本上的在线表现，控制历史预测规则的风险，并讨论从中选择输出的方法\scite{cesabianchi2004generalization}

下面给出一个期望风险的简单结论，说明额外条件怎样进入证明。按独立同分布的到达顺序，仅遍历$D_n$一次，从零参数出发使用固定学习率$\eta>0$和原更新规则。用$f_i$表示看到第$i$个样本之前的分类器；它只依赖$D_{i-1}$，其中$D_0$为空。这里$i$计数已访问的样本位置，与只计实际更新的$t$不同。训练后独立地从$\{1,\ldots,n\}$均匀抽取$I$，返回$f_I$，而不是默认返回最后一次更新后的参数。

\begin{theorem}[随机历史分类器的期望风险]
  \label{thm:classification-perceptron-risk}
  设训练样本独立同分布，测试样本独立来自同一分布。假设存在固定的单位向量$\bm{u}^\ast$及常数$R<\infty$、$\gamma>0$，使
  \[
    \|\tilde{\bm{x}}\|_2\leq R,\qquad
    y\langle\bm{u}^\ast,\tilde{\bm{x}}\rangle\geq\gamma
    \quad\text{几乎必然}.
  \]
  对上述单遍历训练和随机输出规则，有
  \begin{equation}
    \label{eq:classification-perceptron-risk}
    \mathbb{E}_{D_n,I}\!\left[R(f_I\mid D_n,I)\right]
      \leq\frac{R^2}{n\gamma^2}.
  \end{equation}
\end{theorem}

\begin{proof}
  令$e_i=\mathbb{1}\{f_i(\bm{x}_i)\ne y_i\}$。由于$f_i$尚未使用第$i$个样本，独立性给出
  \[
    \mathbb{E}[e_i\mid D_{i-1}]
      =R(f_i\mid D_{i-1}).
  \]
  对两边取期望，再对均匀的$I$平均，就有
  \[
    \mathbb{E}_{D_n,I}[R(f_I\mid D_n,I)]
      =\frac1n\sum_{i=1}^{n}\mathbb{E}[e_i].
  \]
  每次预测错误都会触发更新，零得分且预测正确时也可能更新，故$\sum_i e_i\leq M$。分布层面的共同间隔和长度界使定理~\ref{thm:classification-perceptron-mistakes}几乎必然适用于每个训练序列，于是$M\leq R^2/\gamma^2$，代入即得结论。
\end{proof}

这里的$R$和$\gamma$约束整个数据分布，而不只是碰巧可分的一份训练集。式~\eqref{eq:classification-perceptron-risk}控制的是训练抽样与随机选取历史状态共同平均后的风险，不是对某次训练的高概率保证；当上界大于$1$时也没有提供非平凡信息。它不能直接套给最终参数、平均参数或反复遍历后的参数，因为这些输出与上述随机规则不同，多轮训练时当前样本也不再独立于已有参数。投票感知机等方法正是在如何利用训练历史这一环节作出不同选择，其风险需要相应的分析。

\subsubsection{支持多分类}
\label{sec:classification-perceptron-multiclass}

对$C$类问题，把标签改为$y_i\in\{1,\ldots,C\}$，每类维护一组增广参数$\tilde{\bm{w}}_c=(\bm{w}_c^{\top},b_c)^{\top}$，得分和预测为
\[
  s_c(\bm{x})=\langle\tilde{\bm{w}}_c,\tilde{\bm{x}}\rangle,
  \qquad
  \hat y=\argmax_{c\in\{1,\ldots,C\}}s_c(\bm{x}).
\]
平局仍按最小类别索引处理。这个构造把正确类别的得分与每个竞争类别作比较，也可用\term{Kesler构造}（Kesler construction）把这些比较写成一个扩展空间中的线性可分问题。

遇到$\hat y_i\ne y_i$时，只提高正确类别的参数并降低错误获胜类别的参数：
\begin{equation}
  \label{eq:classification-perceptron-multiclass-update}
  \begin{aligned}
    \tilde{\bm{w}}_{y_i}
      &\leftarrow\tilde{\bm{w}}_{y_i}+\eta\tilde{\bm{x}}_i,\\
    \tilde{\bm{w}}_{\hat y_i}
      &\leftarrow\tilde{\bm{w}}_{\hat y_i}-\eta\tilde{\bm{x}}_i .
  \end{aligned}
\end{equation}
其余类别不变，两个类别的偏置也随各自最后一个分量同时更新。若平局后预测错误，同样更新；若已选择正确类别，则该规则不更新。与二分类中所有零得分均更新的约定不同，这里直接按最终预测是否正确判断。

收敛条件也必须改为多类的得分差条件。把各类参数按行堆成$\tilde{\bm{W}}\in\mathbb{R}^{C\times(d+1)}$；假设存在$\|\bm{U}^\ast\|_{\mathrm{F}}=1$和$\gamma_{\mathrm{mc}}>0$，使对所有$i$和$c\ne y_i$都有
\[
  \langle\bm{U}_{y_i,:}^\ast-\bm{U}_{c,:}^\ast,
    \tilde{\bm{x}}_i\rangle\geq\gamma_{\mathrm{mc}} .
\]
这里$\|\cdot\|_{\mathrm{F}}$是所有矩阵元素平方和的平方根，行向量的内积按对应坐标计算。一次更新的方向矩阵仅两行非零，分别为$\tilde{\bm{x}}_i^{\top}$与$-\tilde{\bm{x}}_i^{\top}$，其平方范数至多为$2R^2$。它与$\bm{U}^\ast$的内积至少为$\gamma_{\mathrm{mc}}$；与当前参数的内积则为$s_{y_i}(\bm{x}_i)-s_{\hat y_i}(\bm{x}_i)\leq0$。因此零初始化时完全沿用二分类证明，得到
\[
  M\eta\gamma_{\mathrm{mc}}
    \leq\|\tilde{\bm{W}}^{(M)}\|_{\mathrm{F}}
    \leq\eta R\sqrt{2M},
  \qquad
  M\leq\frac{2R^2}{\gamma_{\mathrm{mc}}^2}.
\]
有限更新保证因而保留下来，但它要求所有正确类与竞争类之间都具有共同的正间隔，而不是只检查任意一对类别。

\subsubsection{小结}
\label{sec:classification-perceptron-summary}

感知机说明了不估计概率分布也可以学习线性边界：只要存在正间隔分隔面，误差驱动更新就能在有限次修正后找到一个可行解。但它没有规定在可行解之间怎样选择，也没有保证不可分数据上停止。围绕这些限制，形成了不同的扩展。

\term{口袋算法}（Pocket algorithm）处理的是“当前权重未必比过去更好”的问题。常用的带历史最优记录的版本额外保存$\tilde{\bm{w}}_{\mathrm{pocket}}$：每次检查候选参数时，计算全训练集零一错误数，仅在严格改进时替换记录；预算用尽后返回记录参数。它可以用于噪声或矛盾标签数据，但有限预算内只是已检查候选中的最好解，不保证全局最优，反复评估全训练集也有额外开销。Gallant在1990年的论文中系统讨论了Pocket及相关学习方法\scite{gallant1990}

若问题在于原空间不能线性分开，可把输入换成特征映射$\bm{\phi}(\bm{x})$。核函数$K(\bm{x},\bm{z})=\langle\bm{\phi}(\bm{x}),\bm{\phi}(\bm{z})\rangle$直接计算映射后的内积，省去显式构造高维坐标，得到\term{核感知机}（kernel perceptron）。为与前文偏置一致，使用增广映射$(\bm{\phi}(\bm{x}),1)$及其核$\tilde K(\bm{x},\bm{z})=K(\bm{x},\bm{z})+1$。零初始化、单位学习率下，参数在特征空间中是更新样本的线性组合，预测可写为
\begin{equation}
  \label{eq:classification-kernel-perceptron}
  f(\bm{x})=\operatorname{sign}
    \left(\sum_{i=1}^{n}\alpha_i y_i
      [K(\bm{x}_i,\bm{x})+1]\right).
\end{equation}
这里$\alpha_i$是样本$i$的更新次数；遇到非正有符号得分时令$\alpha_i\leftarrow\alpha_i+1$。必须选用正定核，即它在任意有限输入集上生成半正定Gram矩阵，才能作上述内积解释；有限更新还要求增广特征空间内存在正间隔，长度界相应变为$\sqrt{K(\bm{x}_i,\bm{x}_i)+1}$的上界。预测需要访问系数非零的样本，数量可能随训练增长。将内积推广到特征空间的思路可追溯至Aizerman、Braverman与Rozonoer在1964年的势函数研究\scite{aizerman1964}

另一类扩展保留多段训练历史。\term{投票感知机}（voted perceptron）保存历次参数$\tilde{\bm{w}}^{(j)}$及其使用时长$q_j$，以时长作为票数：
\[
  f_{\mathrm{vote}}(\bm{x})
    =\operatorname{sign}\left(
      \sum_{j=1}^{J}q_j
        \operatorname{sign}
          (\langle\tilde{\bm{w}}^{(j)},\tilde{\bm{x}}\rangle)
      \right).
\]
其中$J$为保存的参数状态数，$q_j$按该参数持续使用的样本步数计数，而不只数更新次数。\term{平均感知机}（averaged perceptron）则对相同历史平均参数，
\[
  \bar{\tilde{\bm{w}}}
    =\frac{\sum_{j=1}^{J}q_j\tilde{\bm{w}}^{(j)}}
      {\sum_{j=1}^{J}q_j},
  \qquad
  f_{\mathrm{avg}}(\bm{x})
    =\operatorname{sign}
       (\langle\bar{\tilde{\bm{w}}},\tilde{\bm{x}}\rangle).
\]
平均向量可通过累计和维护，预测仍只需一次线性得分；投票则需保留多个状态。先取符号再投票与先平均再取符号一般不等价。Freund与Schapire在1999年的工作中研究了投票法，也比较了平均法；其大间隔泛化结论有明确的抽样和训练设置，不能解释为二者在任意噪声数据上都优于最终参数\scite{freund1999}

若希望每次更新不仅改正符号，还达到指定得分间隔，\term{被动攻击算法}（passive-aggressive algorithm, PA）把“尽量少改参数”和“满足当前样本的单位间隔”合成一个投影问题。用$\bm{v}$表示待求的新增广参数，
\begin{equation}
  \label{eq:classification-pa-projection}
  \begin{aligned}
    \min_{\bm{v}}\quad&
      \frac12\|\bm{v}-\tilde{\bm{w}}\|_2^2\\
    \text{满足}\quad&
      y_i\langle\bm{v},\tilde{\bm{x}}_i\rangle\geq1 .
  \end{aligned}
\end{equation}
令$\ell_i=\max\{0,1-y_i\langle\tilde{\bm{w}},\tilde{\bm{x}}_i\rangle\}$。约束已满足时参数不变；否则，拉格朗日驻点要求$\bm{v}=\tilde{\bm{w}}+\tau_i y_i\tilde{\bm{x}}_i$，将活跃约束代回即得
\[
  \tau_i=\frac{\ell_i}{\|\tilde{\bm{x}}_i\|_2^2},
  \qquad
  \tilde{\bm{w}}\leftarrow
    \tilde{\bm{w}}+\tau_i y_i\tilde{\bm{x}}_i .
\]
因此PA在间隔足够时“被动”，不足时按违约大小调整。增广方式也意味着这里把偏置变化纳入距离惩罚。噪声下硬性满足当前样本可能调整过大，PA-I可用$\tau_i=\min\{\lambda_{\mathrm{PA}},\ell_i/\|\tilde{\bm{x}}_i\|_2^2\}$限制步长，其中$\lambda_{\mathrm{PA}}>0$控制容忍程度。Crammer等人的2006年论文统一分析了这些在线大间隔方法；它们不是一次更新就求得全训练集上的最优分隔面\scite{crammer2006}

如果输出是一整条词性序列等组合对象，可以把多分类的得分比较推广到结构之间。\term{结构化感知机}（structured perceptron）用联合特征$\bm{\phi}(\bm{x},y)$描述输入与候选结构，预测
\[
  \hat y_i=\argmax_{y\in Y(\bm{x}_i)}
    \langle\bm{\theta},\bm{\phi}(\bm{x}_i,y)\rangle ,
\]
其中$Y(\bm{x}_i)$是候选结构集合，$\bm{\theta}$是与联合特征同维的参数。解码错误时进行
\[
  \bm{\theta}\leftarrow\bm{\theta}
    +\eta\bigl[\bm{\phi}(\bm{x}_i,y_i)
      -\bm{\phi}(\bm{x}_i,\hat y_i)\bigr].
\]
需要的类别或结构偏置也作为联合特征中的相应指示分量处理；对所有候选都相同的常数项会在得分差中消去。在特征差有界、正确结构与所有错误结构之间正间隔可分、且能精确解码的条件下，可沿用多类证明。近似搜索则需另行分析，不能直接继承这一保证。Collins在2002年将这种更新与参数平均用于序列标注，提供了词性标注和短语识别等实例，也为比较感知机与条件随机场等结构化模型提供了参照\scite{collins2002}

这些扩展改变的是不同环节：Pocket选择返回哪个历史解，核方法改变输入表示，投票与平均利用训练轨迹，PA改变单步目标，结构化感知机改变输出及解码方式。它们保留了逐样本修正得分的思路，却不会自动消除各自的可分性、计算或泛化限制。理解这种区分，比把所有扩展都视为原始感知机的无条件改进更有助于选择方法。

\subsection{支持向量机（SVM）}
\label{sec:classification-svm}

感知机回答了一个基本问题：如果训练样本能够被某个超平面分开，怎样通过错误驱动的更新找到一个分离器？但分开训练样本并不意味着分界的位置已经合适。同一批可分数据往往允许许多分离超平面，即使从零参数开始，感知机的结果仍会受样本顺序和更新规则影响。一个紧贴训练样本的边界虽然没有训练错误，却可能在输入受到很小扰动时改变预测。要在这些分离器之间作出选择，还需要说明什么样的边界更值得保留。

\term{支持向量机}（support vector machine，SVM）用训练样本到决策边界的距离回答这一问题：在正确分离样本的前提下，尽量增大最近样本到边界的距离。这个选择准则称为最大间隔。它不是对所有噪声和所有分布都最优的保证，而是在给定特征表示与距离度量后，用一个明确的几何量约束分类器。允许部分样本违反间隔要求，可以处理不可分数据；改变计算距离所依赖的特征表示，又可以得到非线性边界。这两个扩展分别改变约束与表示，不应混为同一种改进。

最优分离超平面的早期研究可以追溯到20世纪60年代。现代支持向量分类器的一项关键进展，是Boser、Guyon与Vapnik在1992年把最优间隔训练和核表示结合起来\scite{boser1992margin}这篇论文从分类器容量与泛化之间的关系出发，用最靠近边界的支持样本表示解，并研究了多项式和径向基等非线性分类函数。在高阶多项式情形中，显式列出全部交互特征可能难以承受，而只计算样本间的核函数可以避开这一展开。论文中的字符识别实验也让这种表示方式与具体分类任务联系起来。

1995年Cortes与Vapnik进一步系统发展了容许训练样本违反间隔约束的支持向量网络\scite{cortes1995}这样，无法严格分开的数据也能进入带有惩罚的优化问题，间隔不再以训练样本全部正确为前提。沿着这条路线，模型表达、误差容忍与实际求解逐渐形成分工：核函数决定怎样比较输入，松弛惩罚决定怎样对待间隔违反，后面介绍的分解算法则降低每次优化所需处理的问题规模。几何建模、凸优化和统计分析因此相互联系，但仍分别回答选择什么模型、如何求解以及怎样控制新样本误差。

主线阅读只需抓住“硬间隔选择边界—软间隔容许违反—核方法改变表示—多分类组织类别竞争”这条路径，并理解训练目标与预测得分。对偶、SMO与风险界属于进阶证明和计算细节，初读时可按需要查阅。再生核Hilbert空间构造和Mercer谱展开已放在简介的第\ref{sec:classic-overview-kernel-validity}节，作为核表示的进阶背景。推导所需的凸优化背景可回顾
\BookChapterRef{chap:convex-learning-and-stability}{第二卷“基础、理论和可信性”的“一般学习的可学习性”章}；
本节还会在首次使用处说明拉格朗日对偶、Slater条件、KKT条件和互补松弛，避免把这些名称当作无需解释的前提。

下面从严格线性可分的情形建立模型，推导对偶与求解过程，再分别讨论优化保证和统计保证。有了这一基础，软间隔与核方法就可以理解为对原有问题的局部修改。多分类进一步改变类别之间的比较方式，可以组合二分类器，也可以直接学习多个类别的相对得分。

\subsubsection{模型结构}
\label{sec:classification-svm-model}

沿用二分类标签约定，设训练集为$D_n=\{(\bm{x}_i,y_i)\}_{i=1}^n$，其中$\bm{x}_i\in\mathbb{R}^d$、$y_i\in\{-1,+1\}$，并假设两类均非空。令实值得分为$g(\bm{x})=\bm{w}^{\top}\bm{x}+b$，其中$\bm{w}\in\mathbb{R}^d$、$b\in\mathbb{R}$且$\bm{w}\ne\bm{0}$。决策超平面是$g(\bm{x})=0$，预测规则为$\hat y(\bm{x})=\operatorname{sign}(g(\bm{x}))$；若得分恰为零，约定预测为$+1$。后文会将零得分保守地计入间隔错误，因此这个平局约定不影响风险上界。

仅比较$y_i g(\bm{x}_i)$还不能比较不同超平面。把$(\bm{w},b)$同时乘以任意正数，分类结果不变，这个数值却会随之放大。为区分参数尺度与真实距离，定义样本的\term{函数间隔}（functional margin）为$m_i=y_i g(\bm{x}_i)$，定义其\term{几何间隔}（geometric margin）为
\begin{equation}
  \gamma_i
  =\frac{y_i(\bm{w}^{\top}\bm{x}_i+b)}
         {\lVert\bm{w}\rVert_2},
  \qquad
  \gamma=\min_{1\le i\le n}\gamma_i.
  \label{eq:classification-svm-margin}
\end{equation}
函数间隔保留了得分的尺度；几何间隔除以法向量长度，得到沿法向方向测量的有符号距离。样本分类正确时，这个距离为正。若输入改变为$\bm{x}_i+\bm{\Delta}$，由Cauchy--Schwarz不等式可得
\[
  y_i g(\bm{x}_i+\bm{\Delta})
  \ge y_i g(\bm{x}_i)
      -\lVert\bm{w}\rVert_2\lVert\bm{\Delta}\rVert_2.
\]
因此，$\lVert\bm{\Delta}\rVert_2<\gamma_i$时预测不会翻转。这说明间隔刻画的是当前特征空间中对输入扰动的容忍程度，而不是对任意分布变化的鲁棒性。

严格线性可分意味着存在$(\bm{w},b)$使所有$m_i>0$。由于样本数有限，可以将这组参数除以$\min_i m_i$，得到$\min_i m_i=1$的规范化表示。在该表示下，最近样本到决策超平面的单侧距离为$1/\lVert\bm{w}\rVert_2$；两个平行超平面$g(\bm{x})=+1$与$g(\bm{x})=-1$之间的总宽度则是$2/\lVert\bm{w}\rVert_2$。本节将前者记为$\gamma$，不把单侧间隔与两侧带宽混用。

最大化$\gamma$于是等价于最小化$\lVert\bm{w}\rVert_2$。平方与正比例因子不改变最优解，却使求导更方便，得到\term{硬间隔}（hard margin）SVM的原始问题：
\begin{equation}
  \begin{aligned}
    \min_{\bm{w},b}\quad&
      \frac12\lVert\bm{w}\rVert_2^2\\
    \text{s.t.}\quad&
      y_i(\bm{w}^{\top}\bm{x}_i+b)\ge1,
      \quad i=1,\ldots,n.
  \end{aligned}
  \label{eq:classification-svm-hard-primal}
\end{equation}
这里的“硬”表示每个样本都必须满足间隔约束，不能用其他样本的改善来补偿某个样本的违反。约束写为大于等于$1$并没有放弃规范化：若一个可行解的最小函数间隔严格大于$1$，将参数一起缩小就能降低目标，所以最优解必有紧约束。

\begin{theorem}[硬间隔解的存在性与唯一性]
  \label{thm:classification-svm-hard-unique}
  若有限训练集严格线性可分且两类均非空，则问题\eqref{eq:classification-svm-hard-primal}存在唯一的规范化最优参数$(\bm{w}^{\ast},b^{\ast})$。目标函数只对$\bm{w}$严格凸，偏置的唯一性还需要两类约束共同确定。
\end{theorem}
\begin{proof}
取一个可行解，并把目标限制在不超过该解目标值的水平集上，则$\bm{w}$有界。固定一个正类样本与一个负类样本，它们的约束分别给出$b$的下界与上界；因为$\bm{w}$有界，$b$也有界。该水平集闭且非空，因此连续目标能在其中取得最小值。

若两个最优解的权重不同，其平均仍可行，但平方范数的严格凸性会使平均解的目标更小，产生矛盾。因此所有最优解共享同一个$\bm{w}^{\ast}$。为研究$b$，对给定$\bm{w}$定义两类样本的投影端点
\[
  a_+(\bm{w})=\min_{i:y_i=+1}\bm{w}^{\top}\bm{x}_i,
  \qquad
  a_-(\bm{w})=\max_{i:y_i=-1}\bm{w}^{\top}\bm{x}_i.
\]
可行偏置恰好满足
\[
  1-a_+(\bm{w})\le b\le-1-a_-(\bm{w}),
\]
所以投影间距$a_+(\bm{w})-a_-(\bm{w})$至少为$2$。若在最优$\bm{w}^{\ast}$处它严格大于$2$，令$\theta=2/(a_+-a_-)<1$，取新权重$\theta\bm{w}^{\ast}$和新偏置$-\theta(a_++a_-)/2$，则两类最近样本的函数间隔都为$1$，其余样本仍满足约束，而非零权重的范数严格减小。这与最优性矛盾。故最优投影间距只能等于$2$，可行偏置区间收缩为单点，$b^{\ast}$也唯一。
\end{proof}

这个证明还表明，最优解的两个支持超平面都至少接触一个训练样本。唯一性属于固定数据与固定特征表示下的规范化硬间隔问题，不意味着更换特征、加入噪声或放松约束以后仍有同样的结论。

\begin{example}[单侧间隔与带宽]
  \label{ex:classification-svm-one-dimensional}
  考虑一维教学数据$(x_1,y_1)=(-1,-1)$、$(x_2,y_2)=(1,+1)$、$(x_3,y_3)=(-2,-1)$和$(x_4,y_4)=(2,+1)$。任意落在$-1$与$1$之间的阈值都能正确分类，但只有阈值$0$使两侧最近距离的较小者达到最大。规范化最优解为$w^{\ast}=1$、$b^{\ast}=0$；最近样本的单侧距离为$1$，两条支持超平面在此退化为两个点$-1$与$1$，其间距离为$2$。更远的两个样本不限制这一最优间隔。
\end{example}

\subsubsection{参数学习}
\label{sec:classification-svm-learning}

原始问题可以交给凸二次规划求解器，但直接优化$\bm{w}$不是唯一途径。每个约束究竟在多大程度上限制了最优边界，可以通过与它对应的拉格朗日乘子表示。转到\term{对偶问题}（dual problem），就是改为优化这些约束权重；它既能揭示哪些样本真正参与决策，也能把训练样本的使用方式压缩为样本间内积。这是后面引入核方法的关键，而不是说对偶在所有$n$、$d$组合下都比原始问题便宜。

将约束写为$1-y_i(\bm{w}^{\top}\bm{x}_i+b)\le0$，引入$\alpha_i\ge0$，记$\bm{\alpha}=(\alpha_1,\ldots,\alpha_n)^{\top}$。拉格朗日函数为
\begin{equation}
  \begin{aligned}
    L(\bm{w},b,\bm{\alpha})
    ={}&\frac12\lVert\bm{w}\rVert_2^2\\
      &+\sum_{i=1}^n\alpha_i
        \bigl[1-y_i(\bm{w}^{\top}\bm{x}_i+b)\bigr].
  \end{aligned}
  \label{eq:classification-svm-hard-lagrangian}
\end{equation}
对一个可行$(\bm{w},b)$，$\sup_{\bm{\alpha}\ge\bm{0}}L$等于原始目标；若有约束违反，对应乘子趋于无穷会使上确界为正无穷。因此原始最优值可写为$\inf_{\bm{w},b}\sup_{\bm{\alpha}\ge\bm{0}}L$。反过来，先对$(\bm{w},b)$取下确界再最大化，得到原始最优值的下界。两种次序能否给出相同的值，需要强对偶条件，不能只根据形式交换极值。

这里可以直接核验\term{Slater条件}（Slater's condition）：它要求存在严格满足不等式约束的点。取一个严格分离器，并把参数放大到所有函数间隔都严格大于$1$，就得到所需点。结合目标凸、约束仿射以及上一小节的最优值可达性，可得强对偶与最优乘子的存在性。

现在计算$q(\bm{\alpha})=\inf_{\bm{w},b}L(\bm{w},b,\bm{\alpha})$。由于$b$没有二次项，若$\sum_i\alpha_i y_i\ne0$，沿$b$的某个方向就有$L\to-\infty$。要使$q$有限，必须满足等式约束；关于$\bm{w}$的极小值则由驻点给出：
\begin{equation}
  \bm{w}=\sum_{i=1}^n\alpha_i y_i\bm{x}_i,
  \qquad
  \sum_{i=1}^n\alpha_i y_i=0.
  \label{eq:classification-svm-stationarity}
\end{equation}
令$\bm{v}=\sum_i\alpha_i y_i\bm{x}_i$，代入后涉及权重的两项变为$\frac12\lVert\bm{v}\rVert_2^2-\lVert\bm{v}\rVert_2^2$，而$b$项为零。展开$\lVert\bm{v}\rVert_2^2$，便得到
\begin{equation}
  \begin{aligned}
    \max_{\bm{\alpha}}\quad&
      q(\bm{\alpha})
      =\sum_{i=1}^n\alpha_i
       -\frac12\sum_{i,j=1}^n
          \alpha_i\alpha_j y_i y_j
          \bm{x}_i^{\top}\bm{x}_j\\
    \text{s.t.}\quad&
      \sum_{i=1}^n\alpha_i y_i=0,\qquad
      \alpha_i\ge0,\quad i=1,\ldots,n.
  \end{aligned}
  \label{eq:classification-svm-hard-dual}
\end{equation}
二次项矩阵的元素为$\bm{Q}_{i,j}=y_i y_j\bm{x}_i^{\top}\bm{x}_j$。对任意$\bm{a}\in\mathbb{R}^n$，有$\bm{a}^{\top}\bm{Q}\bm{a}=\lVert\sum_i a_i y_i\bm{x}_i\rVert_2^2\ge0$，所以$q$是凹函数；它未必严格凹，因而对偶乘子未必唯一。

原始解与对偶解的联系可以用\term{Karush--Kuhn--Tucker条件}（Karush--Kuhn--Tucker conditions，KKT条件）完整表达。它们同时检查解是否满足约束、是否达到驻点，以及一个约束是否真的限制了目标的继续改善。在当前凸问题与Slater条件下，这组条件对最优性既必要又充分：
\begin{equation}
  \begin{aligned}
    &m_i^{\ast}=y_i g^{\ast}(\bm{x}_i)\ge1,
      \qquad \alpha_i^{\ast}\ge0,\\
    &\alpha_i^{\ast}(m_i^{\ast}-1)=0,
      \qquad i=1,\ldots,n,\\
    &\bm{w}^{\ast}
      =\sum_{i=1}^n\alpha_i^{\ast}y_i\bm{x}_i,
      \qquad \sum_{i=1}^n\alpha_i^{\ast}y_i=0.
  \end{aligned}
  \label{eq:classification-svm-hard-kkt}
\end{equation}
其中第三个关系是\term{互补松弛}（complementary slackness）：乘子与约束的剩余量不能同时为正。最后一行是关于权重与偏置的驻点条件。

把$\alpha_i^{\ast}>0$的样本称为\term{支持向量}（support vector），并记其索引集为$S=\{i:\alpha_i^{\ast}>0\}$。互补松弛说明，支持向量一定满足$m_i^{\ast}=1$，而$m_i^{\ast}>1$必有$\alpha_i^{\ast}=0$。反向推断却不成立：落在边界上的样本也可能具有零乘子。例如把示例\ref{ex:classification-svm-one-dimensional}中的$x_2=1$重复一次，正类边界点的乘子总和仍可为$1/2$，但这个总量可以全部分配给其中一个副本。超平面没有改变，对偶表示与支持向量索引集却可能改变。

求得$\bm{\alpha}^{\ast}$后，恢复模型只需
\begin{equation}
  \begin{aligned}
    \bm{w}^{\ast}
      &=\sum_{i\in S}\alpha_i^{\ast}y_i\bm{x}_i,\\
    b^{\ast}
      &=y_k-(\bm{w}^{\ast})^{\top}\bm{x}_k,
        \qquad k\in S.
  \end{aligned}
  \label{eq:classification-svm-recover}
\end{equation}
精确最优解下，任意支持向量都会给出同一偏置；数值计算中可以平均多个合适样本给出的估计。只有非零乘子进入预测，这就是SVM的稀疏表示。但稀疏不等于支持向量一定很少，某些数据集可以让大部分甚至全部训练点成为支持向量。

\paragraph{SMO 算法}
\label{sec:classification-svm-smo}

对偶把未知量换成了$n$个乘子，也带来了新的计算负担。若存储完整内积矩阵，需要$O(n^2)$空间；通用稠密求解中的一次矩阵分解就可能需要$O(n^3)$运算，具体总成本还取决于求解器与迭代次数。为了避开每轮求解大规模二次规划，Platt提出了\term{序列最小优化}（sequential minimal optimization，SMO）：每次只优化一个最小的可动变量组，把整个问题分解为可解析处理的二维子问题\scite{platt1998smo}

为什么不能每次只动一个乘子？因为$\sum_i\alpha_i y_i=0$且$y_i\ne0$，单独改变一个$\alpha_i$会破坏等式约束。选择两个不同索引$i,j$，固定其余乘子，记更新前数值为$\alpha_i^{(t)},\alpha_j^{(t)}$，并令$s=y_i y_j\in\{-1,+1\}$。如果把新的$\alpha_j$写为$u$，等式约束要求
\begin{equation}
  \alpha_i(u)
  =\alpha_i^{(t)}
    +s\bigl(\alpha_j^{(t)}-u\bigr).
  \label{eq:classification-svm-smo-coupling}
\end{equation}
这样二维问题实际上只剩一个自由变量。对硬间隔，要求$u\ge0$且$\alpha_i(u)\ge0$，可行区间$[L,H]$为
\begin{equation}
  \begin{aligned}
    s=+1:\quad&
      L=0,\qquad H=\alpha_i^{(t)}+\alpha_j^{(t)},\\
    s=-1:\quad&
      L=\max\{0,\alpha_j^{(t)}-\alpha_i^{(t)}\},
      \qquad H=+\infty.
  \end{aligned}
  \label{eq:classification-svm-smo-hard-interval}
\end{equation}
同号时两乘子之和固定，因而有有限上界；异号时差固定，二者可以同时增大。这两种区间的差别来自标签符号对等式约束的影响。

为求沿这个区间的最优点，定义$\bm{K}_{i,j}=\bm{x}_i^{\top}\bm{x}_j$，并使用当前实值得分计算误差缓存
\[
  E_k=g^{(t)}(\bm{x}_k)-y_k,\qquad
  g^{(t)}(\bm{x})
  =\sum_{\ell=1}^n\alpha_\ell^{(t)}
     y_\ell\bm{x}_\ell^{\top}\bm{x}+b^{(t)}.
\]
$E_k$保留实值得分相对于标签的偏差，因而包含更新所需的距离信息；经过$\operatorname{sign}$得到的离散预测标签不再包含这部分信息。令$\Delta=u-\alpha_j^{(t)}$。由于$\Delta\alpha_i=-s\Delta$，权重变化为$y_j\Delta(\bm{x}_j-\bm{x}_i)$，代入$q$可以逐项得到
\begin{equation}
  \begin{aligned}
    q(u)-q(\alpha_j^{(t)})
      &=y_j(E_i-E_j)\Delta-\frac12\eta\Delta^2,\\
    \eta
      &=\bm{K}_{i,i}+\bm{K}_{j,j}-2\bm{K}_{i,j}\\
      &=\lVert\bm{x}_i-\bm{x}_j\rVert_2^2\ge0.
  \end{aligned}
  \label{eq:classification-svm-smo-gain}
\end{equation}
这里$q(u)$表示其余乘子固定、$\alpha_i$按式\eqref{eq:classification-svm-smo-coupling}联动后的目标。偏置不出现在对偶目标中；误差差$E_i-E_j$中的$b^{(t)}$也恰好相消。

当$\eta>0$时，一元凹二次函数的无约束最优点与区间最优点分别为
\begin{equation}
  \begin{aligned}
    u_{\mathrm{unc}}
      &=\alpha_j^{(t)}+\frac{y_j(E_i-E_j)}{\eta},\\
    \alpha_j^{(t+1)}
      &=\min\{H,\max\{L,u_{\mathrm{unc}}\}\},\\
    \alpha_i^{(t+1)}
      &=\alpha_i^{(t)}
        +s\bigl(\alpha_j^{(t)}-\alpha_j^{(t+1)}\bigr).
  \end{aligned}
  \label{eq:classification-svm-smo-update}
\end{equation}
剪裁只对一个自由变量进行，再由等式回代另一个；独立剪裁两个乘子会破坏耦合约束。若$L=H$，该变量对没有可行移动，应另选工作集。

$\eta=0$时，式\eqref{eq:classification-svm-smo-gain}变为线性函数，在有限区间上比较两个端点的目标即可；若目标恒定，保留旧值避免无效移动。严格可分的线性硬间隔问题中，$\eta=0$意味着两输入重合，它们不可能异号；同号重复点的目标沿可行线恒定。若出现零曲率、无穷上界且目标沿该方向无限增加，则应检查可分性假设、核矩阵与数值误差，并停止该次更新。

乘子更新以后还要调整偏置。记$\Delta\alpha_i=\alpha_i^{(t+1)}-\alpha_i^{(t)}$，$\Delta\alpha_j$同理。分别要求更新后的第$i$或第$j$个样本具有函数间隔$1$，解得
\begin{equation}
  \begin{aligned}
    b_i={}&b^{(t)}-E_i
       -y_i\Delta\alpha_i\bm{K}_{i,i}
       -y_j\Delta\alpha_j\bm{K}_{i,j},\\
    b_j={}&b^{(t)}-E_j
       -y_i\Delta\alpha_i\bm{K}_{i,j}
       -y_j\Delta\alpha_j\bm{K}_{j,j}.
  \end{aligned}
  \label{eq:classification-svm-smo-bias}
\end{equation}
硬间隔中，若更新后的$\alpha_i>0$，使用$b_i$；否则若$\alpha_j>0$，使用$b_j$。二者都为正且子问题精确求解时，这两个值一致；若都不能用于确定偏置，可暂取$(b_i+b_j)/2$。这一步利用当前工作集的条件更新偏置，其他样本是否满足KKT仍由后续全局扫描检查。

令$\Delta b=b^{(t+1)}-b^{(t)}$，其他样本的缓存可由
\[
  E_k^{(t+1)}
  =E_k^{(t)}
   +y_i\Delta\alpha_i\bm{K}_{i,k}
   +y_j\Delta\alpha_j\bm{K}_{j,k}
   +\Delta b
\]
更新，不必重新计算整个预测和。在示例\ref{ex:classification-svm-one-dimensional}中，从$\bm{\alpha}^{(0)}=\bm{0}$、$b^{(0)}=0$开始，选择$x_1=-1$与$x_2=1$，有$E_1=1$、$E_2=-1$、$\eta=4$，于是得到$\alpha_1=\alpha_2=1/2$和$b=0$。此时$w=1$，其余两个样本已经在间隔外；这一例子恰好一步达到最优，不能据此推断一般数据的迭代次数。

\begin{algorithm}[硬间隔SMO的训练框架]
  \label{alg:classification-svm-smo}
  \begin{algorithmic}[1]
    \Require 两类非空且严格可分的$D_n$，容差$\varepsilon_{\mathrm{opt}}>0$，迭代上限$t_{\max}$
    \Ensure 可行乘子、偏置及停止状态
    \State 初始化$\bm{\alpha}^{(0)}=\bm{0}$、$b^{(0)}=0$、$E_k=-y_k$
    \For{$t=0,\ldots,t_{\max}-1$}
      \State 扫描KKT残差；若全部低于容差，返回“容差满足”
      \State 选取违反KKT的索引$i$，再选有可行改进的$j$
      \State 按式\eqref{eq:classification-svm-smo-hard-interval}计算$L,H$
      \State 求解一元子问题；退化时比较端点或更换变量对
      \State 联动更新两个乘子，更新$b$与误差缓存
    \EndFor
    \State 返回当前参数及“达到迭代上限，尚未认证收敛”
  \end{algorithmic}
\end{algorithm}

对于硬间隔，$\alpha_i=0$时应检查$m_i\ge1$，$\alpha_i>0$时应检查$m_i=1$，实际实现还要给“零乘子”设置数值容差。第二个变量可以优先选择误差差较大或预测目标增益较大的样本，但必须有回退扫描，不能永远忽略某些违反条件的变量。算法\ref{alg:classification-svm-smo}给出的是共同框架，具体选点规则的收敛条件留在下一小节讨论。

SMO及其变种无需存放并分解整个稠密二次规划矩阵，不过核缓存仍可能占用可观空间。LIBSVM采用的是带工作集选择策略的SMO类分解方法，而非原始Platt算法的逐行实现；scikit-learn的\code{SVC}使用LIBSVM作为后端\scite{chang2011libsvm}

在线性问题中，也可以直接维护$\bm{w}$，采用原始问题的次梯度方法或合适的坐标下降算法。LIBLINEAR面向大规模稀疏线性分类提供了专门求解路线，避免把所有训练问题都转成稠密核矩阵运算\scite{fan2008liblinear}

\subsubsection{理论分析}
\label{sec:classification-svm-analysis}

一个训练程序给出低目标值，并不能直接说明它找到了最优解；即便已经找到最优解，也不能据此断言测试误差很低。这两个缺口分别需要优化分析与统计分析来弥补。硬间隔的几何结构让两种分析可以共享范数与间隔，但它们的量词、假设和结论仍须分开。

\paragraph{收敛性分析}
\label{sec:classification-svm-convergence}

凸性消除了非全局的局部极小值：一个满足全部KKT条件的可行解就是全局最优解。这是对解的判别，不是对任意迭代过程的保证。算法仍可能反复选择没有有效移动的变量对，或者一直忽略另一组违反KKT的变量；只知道目标不下降，也无法排除它收敛到一个非最优值。

对SMO而言，式\eqref{eq:classification-svm-smo-gain}说明，精确最大化当前子问题不会降低对偶目标。如果$\eta>0$、无约束最优点可行且更新非零，那么增益为$\eta\Delta^2/2>0$。但这些正增益可以越来越小，单调且有上界只足以推出目标值序列有极限，不能推出有限次更新恰好命中最优解。乘子本身不唯一时，目标值收敛与参数序列收敛还可能是不同问题。

Osuna等对SVM分解求解提出了早期训练方法与分析，关键是怎样通过较小的工作集持续修复违反最优性条件的变量，而不是对任意坐标选择作出无条件保证\scite{osuna1997training}

广义SMO的收敛研究进一步明确了算法规则与收敛结论之间的关系；工作集选择、退化情形处理和子问题求解方式都属于定理的条件\scite{keerthi2002convergence}

实际求解器通常报告“在指定容差内满足KKT”，而不是“有限步得到精确实数最优解”。例如Fan等利用二阶信息选择工作集，并对所给算法建立收敛性质，这类结果应与对应选点规则一起使用\scite{fan2005working}若迭代上限先于容差条件触发，则只能报告当前近似解。对于硬间隔，还必须先有可分性使原始问题可行；软间隔通过松弛变量消除这一可行性障碍，但不会自动替任意求解器建立收敛证明。

\paragraph{泛化分析}
\label{sec:classification-svm-generalization}

统计分析需要把训练中获得的间隔与独立测试样本的风险联系起来。以下分别从函数类的复杂度和删除训练样本后的解变化入手，再与感知机的错误更新界比较。

\label{sec:classification-svm-margin-generalization}

\textbf{基于间隔的泛化界。}最大间隔为什么可能有利于新样本？可以预先限制得分函数的范数，再研究其中具有较大训练间隔的函数。线性分类器的几何边界由符号决定；若不固定得分尺度，单靠$\lVert\bm{w}\rVert_2$小还不足以限制分类能力，因为任意非零参数都可以同比缩小而保持分类不变。范数与函数间隔必须共同出现。

\BookSectionRef{sec:rademacher-complexity}{第二卷“基础、理论和可信性”的“二分类学习的可学习性”章“估计期望风险”节}
已经介绍用Rademacher复杂度控制这类得分函数，
\BookEquationRef{eq:rademacher-margin-risk}{第二卷“基础、理论和可信性”的“二分类学习的可学习性”章中的间隔风险界}
给出了无偏置线性模型的间隔风险界。这里补入SVM常用的偏置项，并明确所有上界都在抽样前固定。设独立同分布样本满足$\lVert\bm{x}\rVert_2\le R$几乎处处，给定$B>0$与$B_b\ge0$，考虑
\[
  \mathcal{F}_{B,B_b}
  =\left\{
    \bm{x}\mapsto\bm{w}^{\top}\bm{x}+b:
    \lVert\bm{w}\rVert_2\le B,\ |b|\le B_b
   \right\}.
\]
这里的输入界约束整个数据分布，而不只是已经看到的训练点。偏置也必须得到处理：若直接让$b$任意大，下面对实值得分函数的复杂度估计便不成立。

为了把得分转换成误差，固定$\rho>0$并使用\term{截断间隔损失}（ramp loss）
\begin{equation}
  \ell_\rho(z)=
  \begin{cases}
    1, & z\le0,\\
    1-z/\rho, & 0<z<\rho,\\
    0, & z\ge\rho.
  \end{cases}
  \label{eq:classification-svm-ramp}
\end{equation}
它把错误分类和零得分记为$1$，把达到函数间隔$\rho$的样本记为$0$，中间线性过渡；其Lipschitz常数是$1/\rho$。它用于风险分析，不是把硬间隔训练目标改成这个非凸损失。

\begin{theorem}[含偏置的间隔风险界]
  \label{thm:classification-svm-margin-bound}
  在上述输入有界假设下，固定$B,B_b,\rho$与$\delta\in(0,1)$。以至少$1-\delta$的概率，对所有$g\in\mathcal{F}_{B,B_b}$同时有
  \begin{equation}
    \begin{aligned}
      R_{01}(g)
      &\le \widehat R_\rho(g)
       +\frac{2(BR+B_b)}{\rho\sqrt n}
       +\sqrt{\frac{\log(1/\delta)}{2n}},\\
      \widehat R_\rho(g)
      &=\frac1n\sum_{i=1}^n\ell_\rho(y_i g(\bm{x}_i)),
    \end{aligned}
    \label{eq:classification-svm-margin-bound}
  \end{equation}
  其中$R_{01}(g)=\mathbb{P}(\hat y(\bm{x})\ne y)$是零一风险。
\end{theorem}
\begin{proof}
因为$\mathbb{1}\{\hat y(\bm{x})\ne y\}\le\ell_\rho(y g(\bm{x}))$，只需控制右侧损失的期望。令$\sigma_1,\ldots,\sigma_n$为独立等概率取$\pm1$的Rademacher变量。固定样本后，有符号得分类的经验复杂度满足
\[
  \begin{aligned}
    &\frac1n\mathbb{E}_{\bm{\sigma}}
      \sup_{\lVert\bm{w}\rVert_2\le B,\ |b|\le B_b}
      \sum_{i=1}^n\sigma_i y_i
       (\bm{w}^{\top}\bm{x}_i+b)\\
    &\quad\le
      \frac{B}{n}\mathbb{E}_{\bm{\sigma}}
        \left\lVert\sum_i\sigma_i y_i\bm{x}_i\right\rVert_2
      +\frac{B_b}{n}\mathbb{E}_{\bm{\sigma}}
        \left|\sum_i\sigma_i y_i\right|\\
    &\quad\le
      \frac{B}{n}\sqrt{\sum_i\lVert\bm{x}_i\rVert_2^2}
      +\frac{B_b}{\sqrt n}
      \le\frac{BR+B_b}{\sqrt n}.
  \end{aligned}
\]
第二步先用Jensen不等式，再展开平方；独立随机符号使交叉项的期望为零。这个上界对所有样本成立，取样本期望后也约束期望Rademacher复杂度。由收缩引理，复合损失类的复杂度至多为该值的$1/\rho$倍；需要零点条件时，可以先减去常数$\ell_\rho(0)$，不改变复杂度。对取值在$[0,1]$内的固定损失类，期望Rademacher风险界给出两倍复杂度项与$\sqrt{\log(1/\delta)/(2n)}$的置信项，即得结论。
\end{proof}

这一推导沿用基于Rademacher复杂度的一致风险分析，关键在于它同时覆盖一个预先给定函数类中的所有候选者，因而也覆盖训练后从中选出的模型\scite{bartlettmendelson2002}

硬间隔规范化解的训练函数间隔至少为$1$，取$\rho=1$时经验间隔损失为零；在该解属于预设函数类的前提下，较小的范数会带来较小的复杂度上界。无偏置时，$\rho/\lVert\bm{w}\rVert_2$正是对应的几何间隔，主要复杂度因子可以理解为“输入半径除以间隔”，而不必直接依赖$d$。

SVM通常不正则化$b$，因此式\eqref{eq:classification-svm-margin-bound}需要单独控制偏置。在两类非空的可分数据上，若已有$\lVert\bm{w}^{\ast}\rVert_2\le B$和输入界$R$，前面的偏置区间证明可给出一个确定的有限偏置界，例如$|b^{\ast}|\le BR+1$。另一种处理是把常数特征增广进输入并同时约束其系数，得到一个不同的范数约束模型。

增广表示也便于说明半径与间隔怎样限制容量。令$\widetilde{\bm{x}}=(\bm{x}^{\top},1)^{\top}\in\mathbb{R}^{d+1}$，在所考察输入域上有$\lVert\widetilde{\bm{x}}\rVert_2\le\widetilde R$，并预先固定$\widetilde\gamma>0$。如果$m$个点的每一种$\pm1$标记都能由某个单位增广参数$\bm{u}\in\mathbb{R}^{d+1}$实现，且每个点均满足$y_i\bm{u}^{\top}\widetilde{\bm{x}}_i\ge\widetilde\gamma$，就称这些点能被以$\widetilde\gamma$的\term{间隔打散}（margin shattering）。记这种可打散点数的最大值为$h_{\widetilde\gamma}$，则
\begin{equation}
  h_{\widetilde\gamma}
  \le\min\left\{
      d+1,\left\lfloor
        \frac{\widetilde R^2}{\widetilde\gamma^2}
      \right\rfloor
    \right\}.
  \label{eq:classification-svm-radius-margin-capacity}
\end{equation}
半径项可以用随机符号证明。对每一组独立随机符号$\sigma_1,\ldots,\sigma_m$，按间隔打散的定义选择对应的单位参数$\bm{u}_{\bm{\sigma}}$，则
\[
  m\widetilde\gamma
  \le\bm{u}_{\bm{\sigma}}^{\top}
       \sum_{i=1}^m\sigma_i\widetilde{\bm{x}}_i
  \le\left\lVert
       \sum_{i=1}^m\sigma_i\widetilde{\bm{x}}_i
     \right\rVert_2.
\]
两边平方并对随机符号取期望，交叉项消失，得到$m^2\widetilde\gamma^2\le\sum_i\lVert\widetilde{\bm{x}}_i\rVert_2^2\le m\widetilde R^2$，所以$m\le\lfloor\widetilde R^2/\widetilde\gamma^2\rfloor$。维数项来自线性相关性：若$m>d+1$，存在非零系数满足$\sum_i a_i\widetilde{\bm{x}}_i=\bm{0}$；按非零$a_i$的符号指定标签，间隔条件却会推出$\bm{u}^{\top}\sum_i a_i\widetilde{\bm{x}}_i\ge\widetilde\gamma\sum_i|a_i|>0$，矛盾。

这里控制的是每一种标签都要保留规定距离的打散能力；普通VC维只要求标签可实现，不要求固定的正间隔。因此，观察到某一份已标记数据具有大间隔，并不会改变整个半空间族的VC维。增广表示中的归一化间隔为$y_i(\bm{w}^{\top}\bm{x}_i+b)/\sqrt{\lVert\bm{w}\rVert_2^2+b^2}$，也不同于标准SVM只用$\lVert\bm{w}\rVert_2$归一化的间隔。式\eqref{eq:classification-svm-radius-margin-capacity}揭示了半径与间隔共同限制容量的原因，具体风险保证仍可使用前面的Rademacher界并明确偏置约定。

如果看过数据后才把$B$设为拟合范数、把$\rho$设为观察到的最小间隔，再直接套用固定参数的概率结论，就遗漏了选择这些量的代价。可以预设一列嵌套范数类并分配失败概率，通过并集界实施\term{结构风险最小化}（structural risk minimization，SRM），也可以用独立验证选择超参数。大间隔理论支持的是这样带有容量约束与选择控制的结论，不是单个训练集上的间隔数值本身就构成泛化保证。

\label{sec:classification-svm-loo}

\textbf{基于支持向量数的留一法界。}间隔界从函数类出发；另一种分析直接考察删除一个样本会不会改变最优解。记$A$为精确求解硬间隔问题的学习器：两类非空且可分时返回唯一超平面，训练集只含一类时返回该类的常值分类器。对$n\ge2$，定义\term{留一法}（leave-one-out，LOO）误差为
\begin{equation}
  \widehat R_{\mathrm{LOO}}(D_n)
  =\frac1n\sum_{i=1}^n
    \mathbb{1}\bigl\{
       A(D_n\setminus\{(\bm{x}_i,y_i)\})(\bm{x}_i)
       \ne y_i
     \bigr\}.
  \label{eq:classification-svm-loo-definition}
\end{equation}
这里删除的是第$i$个观测，即使两个输入相同，也把它们视为不同训练观测。对固定训练集，选定一组最优乘子，仍记$S=\{i:\alpha_i^{\ast}>0\}$。

\begin{theorem}[硬间隔的支持向量计数界]
  \label{thm:classification-svm-loo-bound}
  若$D_n$严格线性可分且两类均非空，则上述学习器满足
  \[
    \widehat R_{\mathrm{LOO}}(D_n)\le\frac{|S|}{n}.
  \]
\end{theorem}
\begin{proof}
考虑一个$\alpha_i^{\ast}=0$的观测。删除它后，原始解仍满足剩余约束；同时删除零乘子，不改变$\bm{w}^{\ast}$的表示和$\sum_j\alpha_j^{\ast}y_j=0$，剩余的互补松弛条件也全部保留。所以原始解与剩余乘子仍满足缩小问题的KKT条件，仍然最优。

删除零乘子观测不会恰好删去某一类的唯一成员：若一类仅有一个观测且其乘子为零，等式约束将迫使另一类的所有非负乘子也为零，从而$\bm{w}^{\ast}=\bm{0}$，与两类硬间隔可行性矛盾。因此缩小后的数据仍有两类，由定理\ref{thm:classification-svm-hard-unique}，其最优超平面与原来完全一致，被删观测仍被正确预测。

留一错误于是只能发生在$S$内，每个这样的观测至多贡献一次错误，求和并除以$n$即得结论。支持向量被删除也未必引起错误，因此这里是上界而非等式。
\end{proof}

这个证明使用了精确最优解与删除零乘子观测后的解不变性，不能原样推广到任意近似求解器、任意软间隔偏置选择或其他不稳定的解选择规则。对偶表示不唯一并不破坏上述论证：固定任意一组最优乘子，删除其中的零系数后，KKT证明仍成立。

若数据独立同分布，且含两类的相关有限样本几乎必然严格线性可分，则每一个被留出的观测都独立于其余$n-1$个观测。再由交换性，留一误差中的各项具有相同期望，得到
\begin{equation}
  \begin{aligned}
    \mathbb{E}_{D_{n-1}}R_{01}(A(D_{n-1}))
      &=\mathbb{E}_{D_n}\widehat R_{\mathrm{LOO}}(D_n)\\
      &\le\frac{\mathbb{E}_{D_n}|S(D_n)|}{n}.
  \end{aligned}
  \label{eq:classification-svm-loo-expectation}
\end{equation}
只含一类的样本采用前述常值规则，并把其$S$记为空集，相同关系仍适用。左侧是用$n-1$个观测训练的学习器的期望风险，不是当前$n$样本模型的条件风险。当前训练集的$|S|/n$可以作为留一误差的可计算上界，但没有额外概率分析时，不能把它宣称为当前模型真实风险的无条件上界。

\label{sec:classification-svm-perceptron-comparison}

\textbf{与感知机Novikoff定理的对比。}感知机的Novikoff错误次数界也包含$(R/\gamma)^2$。在相同输入范数、偏置处理和间隔约定下，这种形式相近来自同一几何事实：较小的输入半径与较大的分离间隔，限制了错误驱动更新能够持续发生的程度。但Novikoff定理计数的是感知机沿训练序列犯错并更新的次数，不是对任意测试分布的风险界。

SVM的间隔风险界则需要独立同分布抽样，并将训练间隔、候选函数类的范数与样本量共同联系起来。感知机不保证找到最小范数解或最大间隔解，SVM也不因为优化了最大间隔就对所有问题取得更低测试误差。两种结果提供了不同层面的信息：前者解释一个在线更新过程何时停止犯训练错误，后者解释在额外统计条件下，什么样的间隔结构有助于控制新样本错误。

\subsubsection{支持几乎线性分类（软间隔）}
\label{sec:classification-svm-soft-margin}

硬间隔要求每个观测都能被严格分开。只要两个相同输入带有相反标签，约束就无解；即使仍然可分，一个位置异常的样本也可能把边界推向不适合大多数数据的位置。解决办法是允许样本违反间隔要求，但把违反量计入目标。这就是\term{软间隔}（soft margin）：间隔仍是选择边界的依据，却不再是每个观测都必须满足的硬性条件。

为第$i$个观测引入\term{松弛变量}（slack variable）$\xi_i\ge0$，表示其函数间隔相对于$1$的不足量。令$\bm{\xi}=(\xi_1,\ldots,\xi_n)^{\top}$，惩罚系数记为$\lambda_{\mathrm{s}}>0$，以免与本章的类别数$C$混淆。标准的线性松弛惩罚形式为
\begin{equation}
  \begin{aligned}
    \min_{\bm{w},b,\bm{\xi}}\quad&
      \frac12\lVert\bm{w}\rVert_2^2
       +\lambda_{\mathrm{s}}\sum_{i=1}^n\xi_i\\
    \text{s.t.}\quad&
      y_i(\bm{w}^{\top}\bm{x}_i+b)\ge1-\xi_i,\\
    &\xi_i\ge0,\qquad i=1,\ldots,n.
  \end{aligned}
  \label{eq:classification-svm-soft-primal}
\end{equation}
在这个目标中，$\lambda_{\mathrm{s}}$越大，一单位间隔违反越昂贵；它惩罚的是间隔不足，而不只是错分次数。固定$(\bm{w},b)$后，$\xi_i$越小目标越低，因而最优松弛量必为$\max\{0,1-y_i g(\bm{x}_i)\}$。

当$\xi_i=0$时，样本位于规范间隔边界或其外侧；当$0<\xi_i<1$时，样本仍被正确分类，但位于间隔内部；$\xi_i=1$表示零得分；$\xi_i>1$才表示落在错误一侧。图\ref{fig:classification-svm-margin}固定$g(\bm{x})=x_1-2$，用同一个边界比较这几种情形。点A与点B都是正类，分别有$y_Ag(\bm{x}_A)=0.5$和$y_Bg(\bm{x}_B)=-0.5$，所以最小松弛量分别为$0.5$与$1.5$。松弛量是在得分尺度上测量的；当$\bm{w}\ne\bm{0}$时，除以$\lVert\bm{w}\rVert_2$才是沿法向方向补足至相应规范间隔边界的距离。

\begin{figure}[htbp]
  \centering
  % 教学构造：两个面板均取 g(x)=x_1-2，规范间隔边界为 x_1=1 与 x_1=3。
\begin{tikzpicture}[foundation picture,x=1cm,y=1cm,font=\figurefont\footnotesize]
  \foreach \shift/\title in {0/硬间隔,5.6/软间隔允许违反约束} {
    \begin{scope}[xshift=\shift cm]
      \fill[FigureYellowFill] (1,0) rectangle (3,4);
      \draw[foundation axis] (0,0) -- (4.35,0) node[right] {$x_1$};
      \draw[foundation axis] (0,0) -- (0,4.2) node[above] {$x_2$};
      \draw[FigureStroke,line width=1.1pt] (2,0) -- (2,4);
      \draw[FigureBlue,dashed,line width=1pt] (1,0) -- (1,4);
      \draw[FigureBlue,dashed,line width=1pt] (3,0) -- (3,4);
      \node[foundation heading] at (2,4.8) {\title};
      \node[below,font=\footnotesize] at (1,-0.1) {$g=-1$};
      \node[below,font=\footnotesize] at (2,-0.1) {$g=0$};
      \node[below,font=\footnotesize] at (3,-0.1) {$g=1$};
      \foreach \x/\y in {0.5/2.2,0.7/3.3,1/1.1}
        \fill[FigureBlue] (\x,\y) circle (2.4pt);
      \foreach \x/\y in {3/2.6,3.6/1.9,3.7/3.4}
        \fill[FigureYellow] (\x-0.08,\y-0.08) rectangle ++(0.16,0.16);
    \end{scope}
  }
  \draw[FigureStroke] (1,1.1) circle (4.5pt);
  \draw[FigureStroke] (3,2.6) circle (4.5pt);
  \draw[{Stealth[length=1.5mm]}-{Stealth[length=1.5mm]},FigureStroke]
    (1,3.9) -- node[above,font=\footnotesize] {$2/\lVert\bm{w}\rVert_2$} (3,3.9);
  \begin{scope}[xshift=5.6cm]
    \fill[FigureYellow] (2.5-0.08,0.9-0.08) rectangle ++(0.16,0.16);
    \node[above left] at (2.5,0.9) {A};
    \fill[FigureYellow] (1.5-0.08,2.9-0.08) rectangle ++(0.16,0.16);
    \node[above right] at (1.5,2.9) {B};
    \draw[FigureRed,line width=1.5pt] (2.5,0.9) -- (3,0.9)
      node[midway,below,text=FigureInk,inner sep=1pt] {$\xi_A=0.5$};
    \draw[FigureRed,line width=1.5pt] (1.5,2.9) -- (3,2.9)
      node[midway,above,text=FigureInk,inner sep=1pt] {$\xi_B=1.5$};
  \end{scope}
  \fill[FigureBlue] (2.7,-1) circle (2.4pt);
  \node[right] at (2.85,-1) {负类};
  \fill[FigureYellow] (5.9,-1.08) rectangle ++(0.16,0.16);
  \node[right] at (6.15,-1) {正类};
\end{tikzpicture}

  \caption{硬间隔与软间隔的教学示意，坐标无量纲。蓝圆为负类，赭方为正类；两个面板都给定$\bm{w}=(1,0)^{\top}$、$b=-2$，黑色实线为决策边界$g=0$，蓝色虚线为间隔边界$g=\pm1$。单侧规范间隔为$1/\lVert\bm{w}\rVert_2=1$，两虚线总宽为$2/\lVert\bm{w}\rVert_2=2$；左图圈出间隔上的样本。右图红色线段直接标出松弛量：A位于$x_1=2.5$，正确分类但在间隔内，$\xi_A=0.5$；B位于$x_1=1.5$，已被错分，$\xi_B=1.5$。右图仅说明给定分界面下的最小松弛量，并非优化后的SVM最优解。}
  \label{fig:classification-svm-margin}
\end{figure}

消去$\bm{\xi}$，式\eqref{eq:classification-svm-soft-primal}等价于
\begin{equation}
  \min_{\bm{w},b}\quad
    \frac12\lVert\bm{w}\rVert_2^2
    +\lambda_{\mathrm{s}}\sum_{i=1}^n
       \max\{0,1-y_i g(\bm{x}_i)\}.
  \label{eq:classification-svm-hinge-objective}
\end{equation}
其中$\ell(z)=\max\{0,1-z\}$是\term{合页损失}（hinge loss）：达到目标间隔以后不再奖励更大的得分，间隔不足时线性增加惩罚。它与前面的截断间隔损失不同，错误一侧的损失没有截断，因而保持凸性。这一形式把“最大间隔加违规代价”与“平方范数正则化加经验损失”联系起来。

惩罚参数的数值还依赖损失是求和还是取平均。将式\eqref{eq:classification-svm-hinge-objective}除以$n\lambda_{\mathrm{s}}$，等价目标是平均合页损失加$\lVert\bm{w}\rVert_2^2/(2n\lambda_{\mathrm{s}})$。所以在不同样本量或不同软件约定之间比较超参数时，应先统一归一化方式。较大的$\lambda_{\mathrm{s}}$允许模型付出更大范数来降低训练损失，较小的值更重视范数控制；是否过拟合或欠拟合还取决于数据、表示与验证结果，不能由单个超参数独立决定。

软间隔的乘子上界来自松弛变量的最优性条件。分别为$1-\xi_i-y_i g(\bm{x}_i)\le0$与$-\xi_i\le0$引入$\alpha_i,\mu_i\ge0$，拉格朗日函数为
\[
  \begin{aligned}
    L={}&\frac12\lVert\bm{w}\rVert_2^2
       +\lambda_{\mathrm{s}}\sum_i\xi_i\\
      &+\sum_i\alpha_i
        \bigl[1-\xi_i-y_i(\bm{w}^{\top}\bm{x}_i+b)\bigr]
       -\sum_i\mu_i\xi_i.
  \end{aligned}
\]
关于$\bm{w},b$的驻点条件仍是式\eqref{eq:classification-svm-stationarity}，关于每个$\xi_i$则要求
\[
  \lambda_{\mathrm{s}}-\alpha_i-\mu_i=0.
\]
结合$\mu_i\ge0$，得到$0\le\alpha_i\le\lambda_{\mathrm{s}}$。取$\bm{w}=\bm{0}$、$b=0$和所有$\xi_i>1$可以严格满足两组不等式，因此不论数据是否可分，Slater条件都成立。消去原始变量后有
\begin{equation}
  \begin{aligned}
    \max_{\bm{\alpha}}\quad&
      \sum_{i=1}^n\alpha_i
       -\frac12\sum_{i,j=1}^n
         \alpha_i\alpha_j y_i y_j
         \bm{x}_i^{\top}\bm{x}_j\\
    \text{s.t.}\quad&
      \sum_{i=1}^n\alpha_i y_i=0,\\
    &0\le\alpha_i\le\lambda_{\mathrm{s}},
      \qquad i=1,\ldots,n.
  \end{aligned}
  \label{eq:classification-svm-soft-dual}
\end{equation}
对偶可行域非空且紧，目标连续，所以对偶最优值可以取得。原始形式同样有最优解；两类非空时，限制目标的水平集可控制$\bm{w}$与$\bm{\xi}$，再由两类约束控制$b$。

软间隔的互补松弛关系多了一组：
\begin{equation}
  \begin{aligned}
    &\alpha_i^{\ast}(m_i^{\ast}-1+\xi_i^{\ast})=0,\\
    &(\lambda_{\mathrm{s}}-\alpha_i^{\ast})\xi_i^{\ast}=0,\\
    &m_i^{\ast}\ge1-\xi_i^{\ast},\qquad \xi_i^{\ast}\ge0.
  \end{aligned}
  \label{eq:classification-svm-soft-kkt}
\end{equation}
联合原始与对偶可行性，可以逐项读出
\begin{equation}
  \begin{aligned}
    \alpha_i^{\ast}=0
      &\ \Longrightarrow\
        \xi_i^{\ast}=0,\quad m_i^{\ast}\ge1,\\
    0<\alpha_i^{\ast}<\lambda_{\mathrm{s}}
      &\ \Longrightarrow\
        \xi_i^{\ast}=0,\quad m_i^{\ast}=1,\\
    \alpha_i^{\ast}=\lambda_{\mathrm{s}}
      &\ \Longrightarrow\
        m_i^{\ast}=1-\xi_i^{\ast}\le1.
  \end{aligned}
  \label{eq:classification-svm-soft-cases}
\end{equation}
边界上的样本可以出现在三种乘子状态中：$\alpha_i^{\ast}=0$允许$m_i^{\ast}=1$，$\alpha_i^{\ast}=\lambda_{\mathrm{s}}$也允许$\xi_i^{\ast}=0$且$m_i^{\ast}=1$。反过来，$m_i^{\ast}>1$必有$\alpha_i^{\ast}=0$，$m_i^{\ast}<1$必有$\alpha_i^{\ast}=\lambda_{\mathrm{s}}$，而$m_i^{\ast}=1$时乘子可以是盒约束中的任意值，只要整体KKT系统成立。达到上界的乘子因此可能对应边界点、间隔内的正确样本或错分样本，需要结合$m_i^{\ast}$才能区分。

范数项依然只对$\bm{w}$严格凸，足以保证最优权重唯一，却不能保证偏置唯一。一个简单反例是两个输入都为零，标签分别为$+1$与$-1$。此时$\bm{w}^{\ast}=\bm{0}$，任意$b\in[-1,1]$都使总合页损失为$2$，因此都最优。该情形也说明，软间隔可能得到常值得分；只有$\bm{w}\ne\bm{0}$时，才可以继续用$1/\lVert\bm{w}\rVert_2$解释规范间隔的几何宽度。

若存在$0<\alpha_k^{\ast}<\lambda_{\mathrm{s}}$的\term{自由支持向量}（free support vector），它既参与决策又未碰到乘子上界，可用$m_k^{\ast}=1$恢复偏置。若没有这种样本，就应使用KKT给出的偏置区间。令$h_i=\sum_j\alpha_j^{\ast}y_j\bm{K}_{j,i}$，则$\alpha_i^{\ast}=0$要求$y_i(h_i+b)\ge1$，$\alpha_i^{\ast}=\lambda_{\mathrm{s}}$要求$y_i(h_i+b)\le1$。正标签的前一条件与负标签的后一条件给出下界，另两种组合给出上界；从两组界的交集中选择一个$b$，才是完整的恢复规则。该区间不一定收缩为单点。

SMO的一元目标与误差公式不变，但可行区间必须同时保证两个乘子都位于盒子内。令$s=y_i y_j$，新的$\alpha_j$区间为
\begin{equation}
  \begin{aligned}
    s=+1:\quad
    L&=\max\{0,\alpha_i^{(t)}+\alpha_j^{(t)}
                    -\lambda_{\mathrm{s}}\},\\
    H&=\min\{\lambda_{\mathrm{s}},
                    \alpha_i^{(t)}+\alpha_j^{(t)}\};\\
    s=-1:\quad
    L&=\max\{0,\alpha_j^{(t)}-\alpha_i^{(t)}\},\\
    H&=\min\{\lambda_{\mathrm{s}},
           \lambda_{\mathrm{s}}+\alpha_j^{(t)}-\alpha_i^{(t)}\}.
  \end{aligned}
  \label{eq:classification-svm-smo-soft-interval}
\end{equation}
这是由盒约束与等式约束共同截出的区间，通常小于$[0,\lambda_{\mathrm{s}}]$。使用式\eqref{eq:classification-svm-smo-update}剪裁并回代即可；当$\eta=0$时，这里的区间总是有限，可以直接比较端点目标。偏置仍按式\eqref{eq:classification-svm-smo-bias}计算，但优先使用更新后满足$0<\alpha_i<\lambda_{\mathrm{s}}$或$0<\alpha_j<\lambda_{\mathrm{s}}$的样本；二者都不自由时，可暂取两个候选偏置的平均，停止时再按全体KKT条件核验。

在可分数据上，足够大的有限惩罚就能使软间隔与硬间隔精确对应。取任意硬间隔最优乘子$\bm{\alpha}^{\mathrm{h}}$，若$\lambda_{\mathrm{s}}>\max_i\alpha_i^{\mathrm{h}}$，则硬间隔原始解连同$\bm{\xi}=\bm{0}$、$\bm{\alpha}^{\mathrm{h}}$和$\mu_i=\lambda_{\mathrm{s}}-\alpha_i^{\mathrm{h}}>0$满足软间隔KKT。由强对偶它也是软间隔最优解，而且正的$\mu_i$迫使所有最优松弛量为零。若数据不可分，则无论把惩罚增大到多少，都不能恢复一个本来就不存在的硬间隔可行解。

\subsubsection{支持非线性分类（核方法）}
\label{sec:classification-svm-kernels}

软间隔改变对间隔违反的容忍程度，核化则把同一个间隔目标放到第\ref{sec:classic-overview-kernel-trick}节定义的特征空间中。以下直接使用简介中的正定核$K$，推导分类SVM的对偶、预测和风险界；特征映射、核合法性与常用核见第\ref{sec:classic-overview-kernel-validity}节和第\ref{sec:classic-overview-common-kernels}节。

将$\bm{w}\in\mathbb{R}^d$换成$\bm{w}\in\mathcal{H}$，将平方范数换成$\lVert\bm{w}\rVert_{\mathcal{H}}^2$，软间隔原始问题的推导仍然成立。更具体地，把任意$\bm{w}$分解为训练特征张成空间内的部分与其正交部分，后者不改变任何训练得分，却只会增加范数。因此最优权重可以限制在有限个训练特征的张成空间中；驻点条件进一步给出$\bm{w}^{\ast}=\sum_i\alpha_i^{\ast}y_i\bm{\phi}(\bm{x}_i)$。核化对偶于是为
\begin{equation}
  \begin{aligned}
    \max_{\bm{\alpha}}\quad&
      \sum_{i=1}^n\alpha_i
       -\frac12\sum_{i,j=1}^n
          \alpha_i\alpha_j y_i y_j
          K(\bm{x}_i,\bm{x}_j)\\
    \text{s.t.}\quad&
      \sum_{i=1}^n\alpha_i y_i=0,\\
    &0\le\alpha_i\le\lambda_{\mathrm{s}},
      \qquad i=1,\ldots,n.
  \end{aligned}
  \label{eq:classification-svm-kernel-dual}
\end{equation}
在映射后严格可分的硬间隔情形，去掉乘子上界即可。新输入的得分与预测为
\begin{equation}
  \begin{aligned}
    g^{\ast}(\bm{x})
      &=\sum_{i\in S}\alpha_i^{\ast}y_i
         K(\bm{x}_i,\bm{x})+b^{\ast},\\
    \hat y(\bm{x})&=\operatorname{sign}(g^{\ast}(\bm{x})).
  \end{aligned}
  \label{eq:classification-svm-kernel-prediction}
\end{equation}
边界在$\mathcal{H}$中仍然是线性的，在原输入空间中则可能是曲线或曲面。SMO也只需把$\bm{K}_{i,j}$换成$K(\bm{x}_i,\bm{x}_j)$；此时$\eta=\lVert\bm{\phi}(\bm{x}_i)-\bm{\phi}(\bm{x}_j)\rVert_{\mathcal{H}}^2\ge0$，前面的区间、零曲率处理与偏置更新都继续适用。

核技巧去掉了对显式特征维数的直接依赖，但没有去掉对样本量的依赖。构建完整Gram矩阵需要$O(n^2)$次核调用，之后还要执行优化迭代；两者合起来才是训练成本。若一次核计算耗时$t_K$，单次预测的主要核计算量是$O(|S|t_K)$。因此，无限维特征可以通过有限次内积运算使用，与“训练成本只由核计算决定”是两回事。

统计分析也可以沿用同一结构。若$K(\bm{x},\bm{x})\le\kappa^2$几乎处处，则$\lVert\bm{\phi}(\bm{x})\rVert_{\mathcal{H}}\le\kappa$；将间隔风险界中的$R$换成$\kappa$、将权重范数换成Hilbert范数即可。控制的是函数范数、间隔与偏置，而不是显式坐标个数。核和超参数若由数据选择，仍需计入相应的选择代价。

\subsubsection{支持多分类}
\label{sec:classification-svm-multiclass}

二分类模型只需判断得分的正负；当候选类别超过两个时，还需要规定多个类别怎样竞争。硬间隔、软间隔和核方法都能作为构件，但组合多个二分类器与直接优化多类间隔会得到不同的训练目标。

现在将标签改为$y_i\in\{1,\ldots,C\}$，其中$C\ge3$，类别索引记为$c$。最直接的推广仍然保留二分类求解器，只改变训练任务的组织方式。\term{一对一}（one-vs-one，OvO）为每对类别$c<c'$训练一个分类器$h_{c,c'}$，仅使用标签为$c$或$c'$的样本，共有$C(C-1)/2$个模型。每个模型在这两个类别之间投一票，预测为
\[
  \hat y(\bm{x})
  =\argmax_{c\in\{1,\ldots,C\}}
     \sum_{c_1<c_2}\mathbb{1}\{h_{c_1,c_2}(\bm{x})=c\}.
\]
其中$c_1,c_2\in\{1,\ldots,C\}$，票数并列时采用预先约定的固定规则，例如选择较小类别编号。不同类别对可能作出不相容的局部判断，投票负责把这些判断汇总成一个输出，并不要求它们对应同一个全局概率模型。

\term{一对其余}（one-vs-rest，OvR）为每个类别$c$构造二分类标签：原标签为$c$时取$+1$，其余取$-1$。训练$C$个实值得分函数$g_c$后，使用
\[
  \hat y(\bm{x})=\argmax_{c\in\{1,\ldots,C\}}g_c(\bm{x}),
\]
并固定平局规则。即使所有得分都为负，或多个类别得分为正，这个最大值规则仍能给出类别。它比较的是独立模型的得分，得分尺度与训练中的类别不平衡可能影响比较，因此通常需要统一预处理、损失约定并评估各类别的错误，而不能将这些数值直接当作互相可比的后验概率。

图\ref{fig:classification-svm-multiclass}把四个类别展开成具体训练任务。沿着一列看，一对一中的每个类别只参与三个任务；一对其余中，每个样本都参与四个任务，但只有一个任务把它标为正类。沿着一行看，前者忽略不属于当前类别对的样本，后者把所有其他类别合并为负类。这种训练标签的差异，决定了两个组合器面对的子问题并不相同。

\begin{figure}[htbp]
  \centering
  % 四类教学构造。行是二分类任务，列是样本原类别；
% +1 / -1 是该任务的训练标签，点号表示不参与该任务。
\begin{tikzpicture}[x=1cm,y=1cm,font=\small]
  \foreach \shift/\title in {0/一对一（OvO）,5.5/一对其余（OvR）} {
    \begin{scope}[xshift=\shift cm]
      \node[font=\figurefont\small] at (2.4,4.95) {\title};
      \node[font=\footnotesize] at (0.5,4.0) {任务};
      \node[font=\footnotesize] at (2.9,4.5) {样本原类别};
      \foreach \c in {1,2,3,4}
        \node at ({1.65+0.8*(\c-1)},4) {$\c$};
      \draw[FigureGrid] (0,3.75) -- (4.65,3.75);
    \end{scope}
  }
  \foreach \row/\pair/\a/\b/\c/\d in {
    1/{1,2}/+1/-1/\cdot/\cdot,
    2/{1,3}/+1/\cdot/-1/\cdot,
    3/{1,4}/+1/\cdot/\cdot/-1,
    4/{2,3}/\cdot/+1/-1/\cdot,
    5/{2,4}/\cdot/+1/\cdot/-1,
    6/{3,4}/\cdot/\cdot/+1/-1} {
    \pgfmathsetmacro{\rowY}{3.65-0.48*\row}
    \node at (0.5,\rowY) {$h_{\pair}$};
    \node at (1.65,\rowY) {$\a$};
    \node at (2.45,\rowY) {$\b$};
    \node at (3.25,\rowY) {$\c$};
    \node at (4.05,\rowY) {$\d$};
  }
  \begin{scope}[xshift=5.5cm]
    \foreach \row/\a/\b/\c/\d in {
      1/+1/-1/-1/-1,
      2/-1/+1/-1/-1,
      3/-1/-1/+1/-1,
      4/-1/-1/-1/+1} {
      \pgfmathsetmacro{\rowY}{3.65-0.48*\row}
      \node at (0.5,\rowY) {$g_{\row}$};
      \node at (1.65,\rowY) {$\a$};
      \node at (2.45,\rowY) {$\b$};
      \node at (3.25,\rowY) {$\c$};
      \node at (4.05,\rowY) {$\d$};
    }
  \end{scope}
  \draw[book arrow] (2.4,0.35) -- (2.4,-0.15);
  \draw[book arrow] (7.9,0.35) -- (7.9,-0.15);
  \node[book node,align=center,minimum width=4.7cm,font=\small] at (2.4,-0.65)
    {6 个模型各投一票\\选择票数最多的类别};
  \node[book node,align=center,minimum width=4.7cm,font=\small] at (7.9,-0.65)
    {比较 4 个模型的得分\\选择得分最大的类别};
\end{tikzpicture}

  \caption{四分类任务中的两种二分类组合方式。各行对应一个待训练模型，各列对应样本的原类别；$+1$和$-1$是为该任务重新编码的标签，点号表示该类样本不参与训练。一对一训练6个模型并汇总类别选票，一对其余训练4个模型并比较实值得分。图中给出训练与预测规则，不预设两种方法会产生相同的分类结果。}
  \label{fig:classification-svm-multiclass}
\end{figure}

OvO模型多，但每个子问题只用两类样本；OvR模型少，但每个子问题都涉及全部样本，并把一个类别与可能很大的其余类别集合比较。所以“模型数较少”不等于“总训练一定更快”。LIBSVM的多类SVM使用OvO，scikit-learn的核分类器\code{SVC}沿用这一训练方式；其输出得分的展示形状与底层实际训练了哪些二分类器是不同层面的设置。类别很多时，还应把模型存储、支持向量重复使用与预测核调用量纳入选择。

如果希望组合器能够容忍若干二分类判断错误，可以为每个类别安排一个彼此分离的编码。\term{纠错输出码}（error-correcting output codes，ECOC）将类别$c$表示成长度为$q$的码字$\bm{M}_{c,:}$，其中码本$\bm{M}\in\{-1,+1\}^{C\times q}$的各行互不相同，每列同时含两种符号。第$j$列定义一个二分类任务，样本$i$的训练标签为$\bm{M}_{y_i,j}$；训练后得到$q$个二分类器$h_j$。预测时形成输出码$\bm{h}(\bm{x})=(h_1(\bm{x}),\ldots,h_q(\bm{x}))$，再选择最近的类别码字：
\begin{equation}
  \begin{aligned}
    \hat y(\bm{x})
      &=\argmin_c d_{\mathrm{H}}
          \bigl(\bm{h}(\bm{x}),\bm{M}_{c,:}\bigr),\\
    d_{\mathrm{H}}(\bm{u},\bm{v})
      &=\sum_{j=1}^q\mathbb{1}\{u_j\ne v_j\}.
  \end{aligned}
  \label{eq:classification-svm-ecoc}
\end{equation}
这里的\term{汉明距离}（Hamming distance）就是两个码字不一致的位置数。要表示$C$个不同类别，至少需要$q\ge\lceil\log_2 C\rceil$；加入冗余位的目的，是增大不同类别之间的距离，而不只是增加分类器个数。Dietterich与Bakiri系统研究了这种将多分类转换为多个编码位学习问题的方法\scite{dietterich1995ecoc}

纠错能力可以直接证明。记码本最小类间距离为$\Delta_{\min}$，若真实类别为$c$，输出码与$\bm{M}_{c,:}$相差$e$位，则对任意$c'\ne c$，三角不等式给出
\[
  d_{\mathrm{H}}(\bm{h},\bm{M}_{c',:})
  \ge d_{\mathrm{H}}(\bm{M}_{c,:},\bm{M}_{c',:})-e
  \ge\Delta_{\min}-e.
\]
只要$2e<\Delta_{\min}$，真实码字就比任何其他码字更近。因此至多$\lfloor(\Delta_{\min}-1)/2\rfloor$位错误可以被最近码字解码纠正。这个结论以实际错误位数为条件；若多个二分类器因为相同特征缺陷而同时犯错，冗余编码本身并不能保证错误位数足够小。

另一条路线不再训练彼此独立的二分类器，而是直接优化类别之间的竞争关系。\term{Crammer--Singer多类SVM}（Crammer--Singer multiclass SVM）要求真实类别得分至少比每个错误类别高一个单位间隔，从而在同一目标中学习所有类别权重。

令$\bm{W}=[\bm{w}_1,\ldots,\bm{w}_C]\in\mathbb{R}^{d\times C}$，使用带偏置的版本$g_c(\bm{x})=\bm{w}_c^{\top}\bm{x}+b_c$，相应优化问题为\scite{crammer2001multiclass}
\begin{equation}
  \begin{aligned}
    \min_{\bm{W},\bm{b},\bm{\xi}}\quad&
      \frac12\sum_{c=1}^C\lVert\bm{w}_c\rVert_2^2
       +\lambda_{\mathrm{s}}\sum_{i=1}^n\xi_i\\
    \text{s.t.}\quad&
      g_{y_i}(\bm{x}_i)-g_c(\bm{x}_i)\ge1-\xi_i,
      \quad c\ne y_i,\\
    &\xi_i\ge0,\qquad i=1,\ldots,n.
  \end{aligned}
  \label{eq:classification-svm-crammer-singer}
\end{equation}
不使用偏置时固定$\bm{b}=\bm{0}$；使用偏置时，可加$\sum_c b_c=0$去掉共同平移的不识别性。消去松弛变量后，每个样本的损失为$\max\{0,\max_{c\ne y_i}[1+g_c(\bm{x}_i)-g_{y_i}(\bm{x}_i)]\}$，即惩罚最强错误竞争者超过允许差距的程度。预测仍取$\argmax_c g_c(\bm{x})$，并可将权重推广到核特征空间。

这一联合目标明确比较了类别之间的相对间隔，OvR则分别比较“该类”与“其余类”。二者的约束数、求解器和数据分解方式不同，训练成本与预测质量也要在相同数据和调参预算下评估。联合建模提供了不同的归纳偏好，并不构成相对于OvO或OvR的普遍性能排序。

\subsubsection{小结}
\label{sec:classification-svm-summary}
\label{sec:classification-svm-assessment}

SVM的主要特点是把分类偏好写进一个明确的优化问题：表示由核决定，复杂度由范数限制，训练误差通过间隔损失进入目标。它不要求预先假设完整的类条件分布，能够用于高维特征，并通过支持向量表示保留决定边界的训练信息。文本分类可以先评估稀疏线性模型，手写字符识别可以比较线性与局部非线性核，生物序列任务则可能需要结构化核。这些应用说明可选择的表示丰富，并不意味着某个核在所有场景都有优势。

计算代价应按任务拆开估计。通用核SVM可能需要二次规模的核存储，优化时间受工作集策略、缓存、条件数和停止精度影响；预测成本又随支持向量数变化。线性SVM可以直接使用稀疏特征与专门求解器，其扩展性不能用稠密核训练的成本概括。是否适合某一数据规模，应根据实际内存、训练时间和预测延迟衡量，而不采用统一的样本量阈值。

输出含义也需要与使用目标一致。$g(\bm{x})$是有符号的分类得分，其绝对值不等于后验概率。需要概率时，可以在未用于拟合该得分函数的校准数据上使用\term{Platt缩放}（Platt scaling），拟合形如
\[
  \widehat p(y=+1\mid\bm{x})
  =\frac1{1+\exp(a_{\mathrm{cal}}g(\bm{x})+b_{\mathrm{cal}})}
\]
的映射；也可以采用\term{保序回归}（isotonic regression），在保持得分顺序的条件下学习非参数的单调概率映射。前者形状受限，后者更灵活但在校准样本较少时可能过拟合，二者都需要独立评估概率质量\scite{niculescu2005probabilities}

实践中也可用交叉拟合得到样本外得分来训练校准器。校准增加了一层需要数据学习的模型，并不会使原始SVM得分自动成为可靠概率。

回到开头的边界选择问题，SVM给出的答案是：在约定的特征几何中，同时衡量边界的范数与样本的间隔不足。这个答案能够通过凸优化实现，也能在输入有界、范数受限和正确抽样假设下获得统计解释。真正使用时，还需共同选择特征尺度、核参数与$\lambda_{\mathrm{s}}$，并把超参数选择、概率校准和最终测试分开。最大间隔、稀疏表示与核化分别解决不同问题；明确它们的条件，才能判断模型的表达能力、计算成本和泛化表现是否符合当前任务。


\subsection{模型对比与总结}
\label{sec:classification-geometry-summary}

本小节先比较LDA、感知机与SVM，再介绍协方差收缩、核判别、稀疏贝叶斯和间隔参数化等代表性扩展。三种基本模型都通过线性得分比较类别，但相同的函数形式并不意味着相同的学习问题。生成式LDA先估计类均值、协方差和先验，再由贝叶斯法则决定边界；Fisher判别则寻找类间散度相对于类内散度最大的投影。感知机只在违反分类条件时更新，在可分训练集上找到一个可行解；这个解依赖样本顺序，并没有被指定为最大间隔解。SVM把间隔写进优化目标，通过范数与训练损失之间的权衡选取分类器。表\ref{tab:classification-linear-comparison}将这些差异放在同一组条件下比较。

\begin{table}[htbp]
  \centering
  \small
  \begin{tabularx}{\linewidth}{@{}>{\RaggedRight\arraybackslash}p{18mm}>{\RaggedRight\arraybackslash}X>{\RaggedRight\arraybackslash}X@{}}
    \toprule
    模型 & 边界由什么决定 & 主要条件与代价 \\
    \midrule
    LDA & 类均值、共享协方差与类先验 & 高斯假设支持概率解释；协方差估计需足够样本或正则化。 \\
    感知机 & 错分样本触发的权重更新 & 正间隔可分性支持有限更新界；不可分时需停止规则或变体。 \\
    线性 SVM & 范数与间隔损失的权衡 & 硬间隔要求可分；软间隔需选择惩罚强度。 \\
    核 SVM & 特征空间中的间隔 & 核定义相似性；训练与预测开销受样本数和支持向量数影响。 \\
    \bottomrule
  \end{tabularx}
  \caption{线性与核化分类器的比较。LDA 一行指共享协方差的高斯生成模型；Fisher 投影本身不要求高斯分布。各模型的理论保证分别针对估计、更新次数或风险，不能直接互换。}
  \label{tab:classification-linear-comparison}
\end{table}

如果不同类别的类内变化方向和尺度明显不同，共享协方差就会限制LDA的表达能力。第\ref{sec:classification-lda-model}节已经定义的二次判别分析（QDA）保留类条件高斯假设，却为每个类别估计独立的协方差矩阵$\hat{\bm{\Sigma}}_c$。删去各类共同的常数项后，其判别得分为
\begin{equation}
  \begin{aligned}
    \delta_c(\bm{x})
      ={}&\log\hat{\pi}_c-\frac12\log\det\hat{\bm{\Sigma}}_c\\
        &-\frac12(\bm{x}-\hat{\bm{\mu}}_c)^\top
          \hat{\bm{\Sigma}}_c^{-1}(\bm{x}-\hat{\bm{\mu}}_c).
  \end{aligned}
  \label{eq:classification-qda-score}
\end{equation}
任意两类得分之差含有关于$\bm{x}$的二次项，所以边界一般是二次曲面。当协方差相同时，二次项相消，又得到 LDA。增加自由度也增加了估计负担：共享协方差模型需要$O(d^2+Cd)$个参数，QDA 则需要$O(Cd^2)$个参数；每一类只有少量样本时，独立协方差容易不稳定甚至奇异。第\ref{sec:classification-gda}节将从生成分布的角度统一这两种模型。

在完全共享与完全独立之间，可以让各类借用一部分合并统计量。\term{正则化判别分析}（regularized discriminant analysis, RDA）通过收缩协方差来控制这种取舍。Friedman 的原始方法包含两个调节方向：将类内协方差向合并协方差收缩，再向球形协方差收缩\scite{friedman1989rda}
记$n_c$为第$c$类样本数，$\hat{\bm{\Sigma}}$为以$n$为分母的合并协方差。一种与原始构造相符的写法是
\begin{equation}
  \begin{aligned}
    \bm{A}_c(a)
      &=\frac{(1-a)n_c\hat{\bm{\Sigma}}_c+an\hat{\bm{\Sigma}}}
              {(1-a)n_c+an},\\
    \tilde{\bm{\Sigma}}_c(a,b)
      &=(1-b)\bm{A}_c(a)
        +b\,\frac{\operatorname{tr}\bm{A}_c(a)}{d}\bm{I},
        \qquad a,b\in[0,1].
  \end{aligned}
  \label{eq:classification-rda}
\end{equation}
当$b=0$时，$a=0$与$a=1$分别对应 QDA 与 LDA 的协方差估计；增大$b$会减弱各方向方差差异，并在平均方差为正时帮助获得可逆估计。两个参数应在训练集内部通过验证选择。单纯在 LDA 与 QDA 之间插值，并不能在两者协方差都奇异时自动解决可逆性问题。

Fisher 判别还可以沿核方法的路线扩展。Mika 等提出的\term{核 Fisher 判别}（kernel Fisher discriminant, KFD）在特征空间中重新计算类均值与散度，再最大化投影后的类间与类内散度比\scite{mika1999kernel}
将投影方向写成$\bm{w}=\sum_{i=1}^n a_i\phi(\bm{x}_i)$，即可把目标表达为关于系数$\bm{a}$和核矩阵的广义 Rayleigh 商，并用正则化处理类内散度的奇异性。它与核 SVM 的区别仍在训练准则：KFD 关注整体类统计量，SVM 关注间隔及其违反程度。KFD 的基本解通常没有支持向量式的稀疏性，预测可能需要访问全部训练样本；若要从投影得分得到后验概率，还需额外建立得分分布或校准模型。

除了用间隔约束产生稀疏性，还可以为不同基函数分别学习先验的收缩强度。\term{相关向量机}（relevance vector machine, RVM）采用
\[
  g(\bm{x})=w_0+\sum_{i=1}^n w_i K(\bm{x}_i,\bm{x}),
  \qquad
  p(w_i\mid a_i)=\mathcal{N}(0,a_i^{-1}),
\]
其中$a_i>0$是每个权重各自的先验精度，即先验方差的倒数。分类时再用适当的概率链接函数把$g(\bm{x})$转换为类别概率。RVM通过最大化对权重积分后得到的边缘似然来估计各个精度，部分$a_i$可能趋向无穷，对应权重的分布随之集中到零。在数值实现中，可以按预定的精度或贡献准则剔除相应基函数，使预测只保留一部分训练样本的基函数。有限但很大的精度只造成强收缩，并不自动把权重变成精确的零；这里的稀疏化依赖精度学习及剔除过程。Tipping 系统介绍了这一稀疏贝叶斯学习框架\scite{tipping2001rvm}
这种做法可以输出模型内的预测分布，也不需要 SVM 的间隔惩罚参数，但仍需选择基函数并处理非凸优化。分类似然还带来近似推断的需要。相关向量在一些任务中比支持向量少，是实验观察而非普遍保证；采用贝叶斯形式也不保证概率在失配数据上自动校准。

如果希望调参时直接联系间隔错误与支持向量比例，可以考虑\term{$\nu$-SVM}（nu-support vector machine）。它把间隔阈值$\rho\geq0$也作为变量，在目标中用参数$\nu$奖励更大的阈值，同时惩罚平均松弛量：
\begin{equation}
  \begin{aligned}
    \min_{\bm{w},b,\rho,\bm{\xi}}\quad&
      \frac12\lVert\bm{w}\rVert_2^2-\nu\rho
      +\frac1n\sum_{i=1}^n\xi_i\\
    \text{满足}\quad&
      y_i(\bm{w}^{\top}\phi(\bm{x}_i)+b)\geq\rho-\xi_i,\\
    &\xi_i\geq0,\qquad \rho\geq0.
  \end{aligned}
  \label{eq:classification-nu-svm}
\end{equation}
在最优$\rho>0$时，其对偶乘子满足$0\leq\alpha_i\leq1/n$且$\sum_i\alpha_i=\nu$。因此至少需要$\nu n$个非零乘子，而每个严格违反间隔的样本都占用一个达到$1/n$的乘子，于是间隔错误比例不超过$\nu$，支持向量比例不低于$\nu$。这里的错误包括正确分类但落在间隔内的样本；结论不是测试错误率保证\scite{scholkopf2000nu}

参数范围还受到类别比例约束。由偏置的驻点条件$\sum_i\alpha_i y_i=0$，正负两类的乘子和都为$\nu/2$；若两类样本数分别为$n_+$和$n_-$，盒约束便要求
\[
  0<\nu\leq\nu_{\max}
  :=\frac{2\min(n_+,n_-)}{n}.
\]
这里的上界不只是$\rho>0$解释所需的附加条件。若$\nu>\nu_{\max}$，令$\bm{w}=\bm{0}$、让常数偏置把全部样本判为多数类，并令少数类松弛量等于$2\rho$、多数类松弛量为零，就得到一条满足约束的射线；其目标值为$(\nu_{\max}-\nu)\rho$，随$\rho\to\infty$趋于$-\infty$。因此，虽然约束集合仍非空，优化问题却无界而没有有限最优解。在允许区间内取得有限最优解后，前述比例结论仍以最优$\rho>0$为条件；$\rho=0$时不能直接套用。

这些扩展分别改变了协方差的共享程度、特征空间、稀疏化机制或间隔参数化方式。实际比较时，应固定特征处理与评价口径，再判断增加的自由度能否由现有样本支持、获得的表示能力是否值得相应的计算开销。仅凭边界是直线还是曲线，无法回答这些问题。
